% \iffalse meta-comment
%
% Copyright 2017-2026 by Tibor Tómács
%
% This work may be distributed and/or modified under the
% conditions of the LaTeX Project Public License, either version 1.3
% of this license or (at your option) any later version.
% The latest version of this license is in
%   http://www.latex-project.org/lppl.txt
% and version 1.3 or later is part of all distributions of LaTeX
% version 2005/12/01 or later.
%
% This work has the LPPL maintenance status `maintained'.
% 
% The Current Maintainer of this work is Tibor Tómács.
%
% \fi
%
% \iffalse
%<*driver>
\ProvidesFile{fgruler.dtx}
%</driver>
%<package>\NeedsTeXFormat{LaTeX2e}[2020/10/01]
%<package>\ProvidesPackage{fgruler}[2026/01/11 v1.7 Package for drawing rulers on the foreground or in the text]
%
%<*driver>
\documentclass{ltxdoc}
\OnlyDescription
\AddToHook{begindocument/before}[doc/hyperref]{\hypersetup{pdfstartview=FitH,allcolors=black}}% see in doc.sty and hypdoc.sty
\usepackage[a4paper,margin=25mm,marginparwidth=15mm,marginparsep=5mm,headsep=5mm,headheight=4mm,footskip=10mm]{geometry}
\usepackage{ctex}
\usepackage{listings,fontawesome,stackengine}
\usepackage[type=none]{fgruler}

\setlength{\labelsep}{0pt}

\lstset{
  basicstyle=\ttfamily,
  columns=fullflexible,
  keepspaces=true,
}

\lstnewenvironment{examplelst}{\lstset{
gobble=2,
belowskip=\bigskipamount,
basicstyle=\ttfamily,
backgroundcolor=\color{blue!10},
framerule=0pt,
frame=trbl,
framesep=5pt,
columns=fullflexible,
xleftmargin=5pt,
xrightmargin=5pt,
keepspaces}}{}

\newcommand{\commandinline}{\lstinline[
literate={<}{{$\langle$}}1{>}{{$\rangle$}}1,
delim={[is][\color{green!50!black}\normalfont\itshape]{!}{!}},
basicstyle=\color{blue!80!black}\ttfamily,
columns=fullflexible,
keepspaces]}

\begin{document}
    \DocInput{./fgruler-zh.dtx}
\end{document}
%</driver>
% \fi
%
% \GetFileInfo{fgruler.sty}
% 
% \title{\texorpdfstring{{\bfseries\sffamily fgruler} 宏包\\{\large v1.7 (2026/01/11)}}{fgruler 宏包}}
% \author{Tibor Tómács\\{\normalsize\href{mailto:tomacs.tibor@gmail.com}{\texttt{tomacs.tibor@gmail.com}}}}
% \date{}
% \maketitle
%
% \fgruler*[color=red!80!black]{upperleft}
% \noindent{\rulerparams{}{\scriptsize\sffamily\color{teal}}{blue}{}{}\ruler{taperight}{\textwidth}}
%
% \section{简介}
% |fgruler| 宏包可在每一页（或当前页）的前景层以绝对位置绘制水平和垂直标尺。
% 借此您可以检查页面版式尺寸。
% 您也可以在正文中绘制各种标尺。
%
% \medskip\noindent
% |fgruler| 宏包依赖以下宏包：|kvoptions|、|etoolbox|、|xcolor|、|graphicx|、|eso-pic|。
%
% \section{加载宏包}
% 使用以下命令加载宏包：
% \begin{flushleft}
% \commandinline|\usepackage[!<宏包选项>!]{fgruler}|
% \end{flushleft}
% 或
% \begin{flushleft}
% \commandinline|\usepackage{fgruler}|\\
% \commandinline|\setfgruler{!<宏包选项>!}|
% \end{flushleft}
% |\setfgruler| 命令也可在正文中使用。
% 
% \section{宏包选项}\label{sec:options}
% 默认情况下，|fgruler| 宏包在每一页的前景层绘制一个方形标尺。
% 以下宏包选项用于设置这些标尺的参数。
% \begin{description}
% \item\commandinline|unit=!<单位>!|\\
% 标尺单位。\commandinline|!<单位>!| 可取的值：
% \begin{description}
% \item\commandinline|cm| 公制标尺（厘米）。默认值。
% \item\commandinline|in| 英制标尺（英寸）。
% \end{description}
%
% \item\commandinline|type=!<类型名称>!|\\
% 决定标尺的原点、方向和长度。\commandinline|!<类型名称>!| 可取的值：
% \begin{description}
% \item\commandinline|upperleft| 默认值。方形标尺（原点：纸张左上角；方向：向下和向右；长度：纸张尺寸）。参见示例 \ref{example:upperleft}。
% \item\commandinline|upperright| 方形标尺（原点：纸张右上角；方向：向下和向左；长度：纸张尺寸）。参见示例 \ref{example:upperright}。
% \item\commandinline|lowerleft| 方形标尺（原点：纸张左下角；方向：向上和向右；长度：纸张尺寸）。参见示例 \ref{example:lowerleft}。
% \item\commandinline|lowerright| 方形标尺（原点：纸张右下角；方向：向上和向左；长度：纸张尺寸）。参见示例 \ref{example:lowerright}。
% \item\commandinline|upperleftT| 方形标尺（原点：文本区域左上角；方向：向下和向右；长度：文本区域尺寸）。参见示例 \ref{example:upperleftT}。
% \item\commandinline|upperrightT| 方形标尺（原点：文本区域右上角；方向：向下和向左；长度：文本区域尺寸）。参见示例 \ref{example:upperrightT}。
% \item\commandinline|lowerleftT| 方形标尺（原点：文本区域左下角；方向：向上和向右；长度：文本区域尺寸）。参见示例 \ref{example:lowerleftT}。
% \item\commandinline|lowerrightT| 方形标尺（原点：文本区域右下角；方向：向上和向左；长度：文本区域尺寸）。参见示例 \ref{example:lowerrightT}。
% \item\commandinline|alledges| 在纸张所有边缘都绘制标尺。
% 顶部标尺原点：纸张左上角（方向：向右）。
% 底部标尺原点：纸张左下角（方向：向右）。
% 左侧标尺原点：纸张左上角（方向：向下）。
% 右侧标尺原点：纸张右上角（方向：向下）。长度：纸张尺寸。参见示例 \ref{example:alledges}。
% \item\commandinline|alledges*| 与 |alledges| 选项类似，但底部标尺原点为纸张右下角（方向：向左），
% 左侧标尺原点为纸张左下角（方向：向上）。
% 参见示例 \ref{example:alledges*}。
% \item\commandinline|alledgesT| 与 |alledges| 选项类似，但绘制在文本区域的边缘。参见示例 \ref{example:alledgesT}。
% \item\commandinline|alledgesT*| 与 |alledges*| 选项类似，但绘制在文本区域的边缘。参见示例 \ref{example:alledgesT*}。
% \item\commandinline|user| 每一对 \commandinline|!<单位>!|--\commandinline|!<类型名称>!| 都会激活一个 \commandinline|\fgrulertype{!<单位>!}{!<类型名称>!}| 命令。
% 因此 \commandinline|\fgrulertype{!<单位>!}{user}| 的展开结果决定了此选项的效果。
% 在使用 \commandinline|\fgrulerdefuser{!<代码>!}| 命令后，它将变为 \commandinline|!<代码>!|，
% 其中在 \commandinline|!<代码>!| 中您必须以 \commandinline|\fgrulerunit| 引用单位。
% 例如，在执行\\ |\fgrulerdefuser{\fgrulertype{\fgrulerunit}{alledges}}|\\ 之后，|type=user| 选项将等价于 |type=alledges|。
% 另请参见示例 \ref{example:user-first}--\ref{example:user-last}。
% \item\commandinline|none| 不绘制标尺。
% \end{description}
%
% \item\commandinline|type={!<类型名称列表>!}|\\
% \commandinline|!<类型名称列表>!| 是以逗号分隔的有效类型名称列表。
% 这些类型将被组合使用。
% 例如 |type={alledges,alledgesT}|。
% 结果参见示例 \ref{example:combining}。
%
% \item\commandinline|hshift=!<长度>!|\\
% 当 \commandinline|!<类型名称>!| 为 |upperleft|、|lowerleft|、|upperright| 或 |lowerright| 时，标尺的水平偏移。
% 当 \commandinline|!<类型名称>!| 为 |upperleft| 或 |lowerleft| 时，偏移方向向右。
% 当 \commandinline|!<类型名称>!| 为 |upperright| 或 |lowerright| 时，偏移方向向左。
% 默认值：|hshift=0cm|。
% 参见示例 \ref{example:shift-first}--\ref{example:shift-last}。
%
% \item\commandinline|vshift=!<长度>!|\\
% 当 \commandinline|!<类型名称>!| 为 |upperleft|、|lowerleft|、|upperright| 或 |lowerright| 时，标尺的垂直偏移。
% 当 \commandinline|!<类型名称>!| 为 |upperleft| 或 |upperright| 时，偏移方向向下。
% 当 \commandinline|!<类型名称>!| 为 |lowerleft| 或 |lowerright| 时，偏移方向向上。
% 默认值：|vshift=0cm|。
% 参见示例 \ref{example:shift-first}--\ref{example:shift-last}。
%
% \item\commandinline|color=!<颜色名称>!|\\
% 标尺颜色（参见 |xcolor| 宏包）。默认值：|color=black|。
% 参见示例 \ref{example:showframe-color}。
%
% \item\commandinline|numsep=!<长度>!|\\
% 数字与标尺之间的间距。默认值：|numsep=3pt|。
%
% \item\commandinline|markthick=!<长度>!|\\
% 刻度线粗细。默认值：|markthick=0.4pt|。
%
% \item\commandinline|marklength=!<长度>!|\\
% 整数单位处的刻度线长度（参见红色刻度）：
% {\fgrulernoborderline\fgrulercaptioncm{}\fgrulerthickcm{}{}{1pt}\fgrulercolorcm{}{}{red}\ruler{rightup}{3cm}}
% 默认值：|marklength=2mm|。其他刻度的长度参见第 \ref{sec:additional} 节。
%
% \item\commandinline|numfont=!<字体类型>!|\\
% 数字字体类型。默认值：|numfont=\scriptsize\sffamily|。此选项仅可用于 |\setfgruler| 和 |\fgruler*| 命令（参见第 \ref{sec:currentpage} 节）。
%
% \item\commandinline|showframe| 或 \commandinline|showframe=true|\\
% 绘制可见的文本区域和边注区域框线，以及页眉和页脚线。
% 其颜色和粗细由 |color| 和 |markthick| 选项决定。
% 参见示例 \ref{example:showframe-color}。
%
% \item\commandinline|showframe=false|\\
% 停用 |showframe| 选项。
%
% \item\commandinline|nonefgrulers|\\
% 终止前景层上所有标尺的绘制，包括由 |\fgruler| 或 |\fgruler*| 生成的标尺（参见第 \ref{sec:currentpage} 节）。
% 但由 |\ruler| 和 |\squareruler| 绘制的标尺（参见第 \ref{sec:intext} 节）不会消失。
% 此外它还会停用 |showframe| 选项。
% 在此情况下，|fgruler| 宏包不会加载 |eso-pic| 宏包。
% 此选项仅在导言区有效。
%
% 建议在以下两种情况下使用：
% \begin{itemize}
% \item 仅在正文中绘制标尺，不需要检查功能时。
% \item 临时暂停检查功能时。
% \end{itemize}
%
% |type=none| 与 |nonefgrulers| 选项并不相同。区别在于：
% \begin{itemize}
% \item |type=none| 不会终止 |\fgruler| 和 |\fgruler*| 命令以及 |showframe| 选项。
% \item |type=none| 可在文档任意位置修改。
% \item 如果使用 |type=none| 选项而不使用 |nonefgrulers|，|fgruler| 宏包仍会加载 |eso-pic| 宏包。
% \end{itemize}
% \end{description}
% 
% \section{在当前页前景层绘制标尺}\label{sec:currentpage}
% \begin{description}
% \item\commandinline|\fgruler[!<单位!>]{!<类型名称>!}{!<水平偏移>!}{!<垂直偏移>!}|\\
% 在本命令展开所在页的前景层绘制一个方形标尺。
% 您可以在同一页使用多个 |\fgruler| 命令。
%
% 宏包选项（参见第 \ref{sec:options} 节）对此命令同样有效，但 |type|、|hshift| 和 |vshift| 除外。
%
% 如果在导言区使用了 |nonefgrulers| 宏包选项，此命令无效。
% \begin{description}
% \item\commandinline|!<单位!>| 选项：
% \commandinline|cm|、\commandinline|in|（参见第 \ref{sec:options} 节）。
% 其默认值与 |unit=|\commandinline|!<单位!>| 宏包选项指定的值相同。
%
% \item\commandinline|!<类型名称!>| 参数：
% \commandinline|upperleft|、\commandinline|upperright|、\commandinline|lowerleft|、\commandinline|lowerright|（参见第 \ref{sec:options} 节）。
%
% \item\commandinline|!<水平偏移>!|
% 水平偏移量。
% 当 \commandinline|!<类型名称>!| 为 |upperleft| 或 |lowerleft| 时偏移方向向右，否则向左。
%
% \item\commandinline|!<垂直偏移>!|
% 垂直偏移量。
% 当 \commandinline|!<类型名称>!| 为 |upperleft| 或 |upperright| 时偏移方向向下，否则向上。
% \end{description}
%
% 示例：|\fgruler[in]{upperright}{1in}{2.5in}|
% \end{description}
% \begin{description}
% \item\commandinline|\fgruler*[!<宏包选项!>]{!<类型名称>!}|\\
% 在本命令展开所在页的前景层绘制一个标尺。
% 您可以在同一页使用多个 |\fgruler*| 命令。
%
% 宏包选项（参见第 \ref{sec:options} 节）对此命令同样有效，但 |type| 除外。
%
% 如果在导言区使用了 |nonefgrulers| 宏包选项，此命令无效。
% \begin{description}
% \item\commandinline|!<宏包选项!>|
% 第 \ref{sec:options} 节中的所有选项，但 |type|、|showframe| 和 |nonefgrulers| 除外。
%
% \item\commandinline|!<类型名称!>|
% 第 \ref{sec:options} 节中 |type| 选项的所有可能值，但 |none| 值除外。
% \end{description}
%
% 示例：|\fgruler*[color=red,type=in]{alledges}|
%
% 请注意以下两条命令等价：\\
% |\fgruler*[unit=in,hshift=1in,vshift=2.5in]{upperright}|\\
% |\fgruler[in]{upperright}{1in}{2.5in}|
% \end{description}
% 
% \section{在正文中绘制标尺}\label{sec:intext}
% \begin{description}
% \item\commandinline|\ruler[!<单位!>]{!<类型名称>!}{!<长度>!}|\\
% 绘制一个水平或垂直标尺。
% 标尺底部与周围文本的基线对齐。
% 宏包选项（参见第 \ref{sec:options} 节）对此命令无效。
% \begin{description}
% \item\commandinline|!<单位!>| 选项：\\
% \commandinline|cm| 公制标尺（厘米）。默认选项。\\
% \commandinline|in| 英制标尺（英寸）。
%
% \item\commandinline|!<类型名称!>| 参数：\\
% \commandinline|downright | 方向：向下。数字在右侧。\\
% \commandinline|downleft  | 方向：向下。数字在左侧。\\
% \commandinline|upright   | 方向：向上。数字在右侧。\\
% \commandinline|upleft    | 方向：向上。数字在左侧。\\
% \commandinline|rightdown | 方向：向右。数字在下侧。\\
% \commandinline|rightup   | 方向：向右。数字在上侧。\\
% \commandinline|leftdown  | 方向：向左。数字在下侧。\\
% \commandinline|leftup    | 方向：向左。数字在上侧。\\
% \commandinline|taperight | 卷尺（方向：向右）。参见示例 \ref{example:tape}。\\
% \commandinline|tapeleft  | 卷尺（方向：向左）。参见示例 \ref{example:tape}。
%
% \item\commandinline|!<长度>!|
% 标尺长度。
% \end{description}
%
% 示例：|\ruler{rightdown}{5cm}| \ruler{rightdown}{5cm}
%
% \item\commandinline|\ruler*[!<单位!>]{!<类型名称>!}{!<长度>!}|\\
% 与 |\ruler| 类似，但标尺顶部与周围文本的基线对齐。
%
% 示例：
% |\ruler*{taperight}{5cm}|
% \ruler*{taperight}{5cm}
%
% \item\commandinline|\squareruler[!<单位!>]{!<类型名称>!}{!<宽度>!}{!<高度>!}|\\
% 绘制一个方形标尺。
% 方形标尺底部与周围文本的基线对齐。
% 宏包选项（参见第 \ref{sec:options} 节）对此命令无效。
% \begin{description}
% \item\commandinline|!<单位!>| 选项：\\
% \commandinline|cm| 公制标尺（厘米）。默认选项。\\
% \commandinline|in| 英制标尺（英寸）。
%
% \item\commandinline|!<类型名称!>| 参数：\\
% \commandinline|upperleft | 方向：向下和向右。\\
% \commandinline|upperright| 方向：向下和向左。\\
% \commandinline|lowerleft | 方向：向上和向右。\\
% \commandinline|lowerright| 方向：向上和向左。
%
% \item\commandinline|!<宽度>!|
% 方形标尺宽度。
%
% \item\commandinline|!<高度>!|
% 方形标尺高度。
% \end{description}
%
% 示例：
% |\squareruler{upperleft}{5cm}{1cm}|
% \squareruler{upperleft}{5cm}{1cm}
%
% \item\commandinline|\squareruler*[!<单位!>]{!<类型名称>!}{!<宽度>!}{!<高度>!}|\\
% 与 |\squareruler| 类似，但方形标尺顶部与周围文本的基线对齐。
%
% 示例：
% |\squareruler*{upperleft}{5cm}{1cm}|
% \squareruler*{upperleft}{5cm}{1cm}
%
% \item\commandinline|\rulerparams{!<刻度粗细>!}{!<数字字体>!}{!<颜色>!}{!<刻度长度>!}{!<数字间距>!}|\\
% 设置由 |\ruler| 或 |\squareruler| 绘制的标尺参数。
% 如果某个参数为空，则该参数不被修改。
% \begin{description}
% \item\commandinline|!<刻度粗细>!| 刻度线粗细。默认值：|0.4pt|
% \item\commandinline|!<数字字体>!| 数字字体类型。默认值：|\scriptsize\sffamily|
% \item\commandinline|!<颜色>!| 标尺线颜色。默认值：|black|
% \item\commandinline|!<刻度长度>!| 整数单位处的刻度线长度。默认值：|2mm|
% \item\commandinline|!<数字间距>!| 数字与标尺之间的间距。默认值：|3pt|
% \end{description}
%
% 例如，|\rulerparams{}{}{red}{}{}| 将标尺颜色改为红色。
%
% \item\commandinline|\rulerparamsfromfg|\\
% 从当前前景层标尺参数设置标尺参数。
%
% \item\commandinline|\rulernorotatenum|\\
% 默认情况下，垂直标尺的数字（由 |\ruler| 或 |\squareruler| 生成）旋转 90$^\circ$。
% 此命令取消该旋转效果。
% 此命令仅可在正文中展开。
%
% 示例：|{\rulernorotatenum\ruler{upright}{1cm}}| {\rulernorotatenum\ruler{upright}{1cm}}
%
% \item\commandinline|\rulerrotatenum|\\
% 在 |\rulernorotatenum| 之后，重新启用数字旋转。
% 此命令仅可在正文中展开。
% \end{description}
% 
% \section{附加设置命令}\label{sec:additional}
% 以下命令对 |fgruler| 宏包绘制的所有标尺均有效。
% \begin{description}
% \item\commandinline|\fgrulerstartnum{!<数字>!}|\\
% \commandinline|!<数字>!| 为非负整数，作为水平和垂直标尺上的起始数字。
% 默认值：|\fgrulerstartnum{0}|
%
% 示例：|{\fgrulerstartnum{5}\squareruler{lowerleft}{3cm}{1cm}}| {\fgrulerstartnum{5}\squareruler{lowerleft}{3cm}{1cm}}
%
% \item\commandinline|\fgrulerstartnumh{!<数字>!}|\\
% \commandinline|!<数字>!| 为非负整数，作为水平标尺上的起始数字。
% 默认值：|\fgrulerstartnumh{0}|
%
% 示例：|{\fgrulerstartnumh{5}\squareruler{lowerleft}{3cm}{1cm}}| {\fgrulerstartnumh{5}\squareruler{lowerleft}{3cm}{1cm}}
%
% \item\commandinline|\fgrulerstartnumv{!<数字>!}|\\
% \commandinline|!<数字>!| 为非负整数，作为垂直标尺上的起始数字。
% 默认值：|\fgrulerstartnumv{0}|
%
% 示例：|{\fgrulerstartnumv{5}\squareruler{lowerleft}{3cm}{1cm}}| {\fgrulerstartnumv{5}\squareruler{lowerleft}{3cm}{1cm}}
%
% \item\commandinline|\fgrulernoborderline|\\
% 默认情况下，标尺的一侧有边框线。
% 此命令可使其消失。
%
% 示例：|{\fgrulernoborderline\ruler{rightup}{3cm}}| {\fgrulernoborderline\ruler{rightup}{3cm}}
%
% \item\commandinline|\fgrulerborderline|\\
% 在 |\fgrulernoborderline| 之后，重新启用先前的默认边框效果。
%
% \item\commandinline|\fgrulercaptioncm{!<标注>!}|\\
% 公制标尺中的单位标注。
% 默认值：|\fgrulercaptioncm{cm}|
%
% 示例：|{\fgrulercaptioncm{}\ruler{rightup}{3cm}}| {\fgrulercaptioncm{}\ruler{rightup}{3cm}}
%
% \item\commandinline|\fgrulercaptionin{!<标注>!}|\\
% 英制标尺中的单位标注。
% 默认值：|\fgrulercaptionin{inch}|
%
% \item\commandinline|\fgrulerdefnum{!<定义>!}|\\
% 标尺数字由 \commandinline|fgrulernum| 计数器决定。
% 其当前值由 \commandinline|\thefgrulernum| 输出。
% 其默认定义为 |\def\thefgrulernum{\arabic{fgrulernum}}|，等价于 |\fgrulerdefnum{\arabic{fgrulernum}}|。
%
% 示例：|{\fgrulerdefnum{}\fgrulercaptioncm{}\ruler{rightdown}{2cm}}| {\fgrulerdefnum{}\fgrulercaptioncm{}\ruler{rightdown}{2cm}}
%
% \item\commandinline|\fgrulerratiocm{!<比值1>!}{!<比值2>!}|\\
% 公制标尺中的刻度长度比值。
% 如果某个参数为空，则该参数不被修改。
% \begin{description}
% \item\commandinline|!<比值1>!|
% $k/10$\,cm 处的刻度长度比值，其中 $k$ 为正整数且不被 5 整除。\\
% {\fgrulerthickcm{1pt}{}{}\fgrulercolorcm{red}{}{}\ruler{rightup}{3cm}}\\
% 例如，若此比值为 0.5，整数单位处刻度长度为 2\,mm，则该刻度长度为 $0.5\cdot 2\,\mathrm{mm}=1\,\mathrm{mm}$。
%
% \item\commandinline|!<比值2>!|
% $k/2$\,cm 处的刻度长度比值，其中 $k$ 为正奇数。\\
% {\fgrulerthickcm{}{1pt}{}\fgrulercolorcm{}{red}{}\ruler{rightup}{3cm}}
% \end{description}
%
% 默认值：|\fgrulerratiocm{0.5}{0.75}|
%
% \item\commandinline|\fgrulerratioin{!<比值1>!}{!<比值2>!}{!<比值3>!}{!<比值4>!}|\\
% 英制标尺中的刻度长度比值。
% 如果某个参数为空，则该参数不被修改。
% \begin{description}
% \item\commandinline|!<比值1>!|
% $k/16$\,inch 处的刻度长度比值，其中 $k$ 为正奇数。\\
% {\fgrulerthickin{1pt}{}{}{}{}\fgrulercolorin{red}{}{}{}{}\ruler[in]{rightup}{2in}}
%
% \item\commandinline|!<比值2>!|
% $k/8$\,inch 处的刻度长度比值，其中 $k$ 为正奇数。\\
% {\fgrulerthickin{}{1pt}{}{}{}\fgrulercolorin{}{red}{}{}{}\ruler[in]{rightup}{2in}}
%
% \item\commandinline|!<比值3>!|
% $k/4$\,inch 处的刻度长度比值，其中 $k$ 为正奇数。\\
% {\fgrulerthickin{}{}{1pt}{}{}\fgrulercolorin{}{}{red}{}{}\ruler[in]{rightup}{2in}}
%
% \item\commandinline|!<比值4>!|
% $k/2$\,inch 处的刻度长度比值，其中 $k$ 为正奇数。\\
% {\fgrulerthickin{}{}{}{1pt}{}\fgrulercolorin{}{}{}{red}{}\ruler[in]{rightup}{2in}}
% \end{description}
%
% 默认值：|\fgrulerratioin{0.25}{0.375}{0.625}{0.75}|
%
% \item\commandinline|\fgrulerthickcm{!<粗细1>!}{!<粗细2>!}{!<粗细3>!}|\\
% 公制标尺中的刻度线粗细。
% 如果某个参数为空，则该参数不被修改。
% \begin{description}
% \item\commandinline|!<粗细1>!|
% $k/10$\,cm 处的刻度线粗细，其中 $k$ 为正整数且不被 5 整除。
%
% \item\commandinline|!<粗细2>!|
% $k/2$\,cm 处的刻度线粗细，其中 $k$ 为正奇数。
%
% \item\commandinline|!<粗细3>!|
% 整数单位处的刻度线粗细。
% \end{description}
%
% 默认值分别由 |\rulerparams| 的 \commandinline|!<刻度粗细>!| 参数和 |markthick| 选项给出。
%
% 示例：\\
% |{\fgrulerthickcm{}{}{2pt}|\\
% |\rulerparams{}{}{}{5mm}{}|\\
% |\fgrulernoborderline|\\
% |\ruler{rightdown}{3cm}}|\\[2mm]
% {\fgrulerthickcm{}{}{2pt}
% \rulerparams{}{}{}{5mm}{}
% \fgrulernoborderline
% \ruler{rightdown}{3cm}}
%
% \item\commandinline|\fgrulerthickin{!<粗细1>!}{!<粗细2>!}{!<粗细3>!}{!<粗细4>!}{!<粗细5>!}|\\
% 英制标尺中的刻度线粗细。
% 如果某个参数为空，则该参数不被修改。
% \begin{description}
% \item\commandinline|!<粗细1>!|
% $k/16$\,inch 处的刻度线粗细，其中 $k$ 为正奇数。
%
% \item\commandinline|!<粗细2>!|
% $k/8$\,inch 处的刻度线粗细，其中 $k$ 为正奇数。
%
% \item\commandinline|!<粗细3>!|
% $k/4$\,inch 处的刻度线粗细，其中 $k$ 为正奇数。
%
% \item\commandinline|!<粗细4>!|
% $k/2$\,inch 处的刻度线粗细，其中 $k$ 为正奇数。
%
% \item\commandinline|!<粗细5>!|
% 整数单位处的刻度线粗细。
% \end{description}
%
% 默认值分别由 |\rulerparams| 的 \commandinline|!<刻度粗细>!| 参数和 |markthick| 选项给出。
%
% 示例：\\
% |{\fgrulerthickin{}{}{}{}{2pt}|\\
% |\rulerparams{}{}{}{5mm}{}|\\
% |\fgrulernoborderline|\\
% |\ruler[in]{rightdown}{3in}}|\\[2mm]
% {\fgrulerthickin{}{}{}{}{2pt}
% \rulerparams{}{}{}{5mm}{}
% \fgrulernoborderline
% \ruler[in]{rightdown}{3in}}
%
% \item\commandinline|\fgrulercolorcm{!<颜色1>!}{!<颜色2>!}{!<颜色3>!}|\\
% 公制标尺中的刻度颜色。
% 如果某个参数为空，则该参数不被修改。
% \begin{description}
% \item\commandinline|!<颜色1>!|
% $k/10$\,cm 处的刻度颜色，其中 $k$ 为正整数且不被 5 整除。
%
% \item\commandinline|!<颜色2>!|
% $k/2$\,cm 处的刻度颜色，其中 $k$ 为正奇数。
%
% \item\commandinline|!<颜色3>!|
% 整数单位处的刻度颜色。
% \end{description}
%
% 默认值分别由 |\rulerparams| 的 \commandinline|!<颜色>!| 参数和 |color| 选项给出。
%
% 示例：\\
% |{\fgrulercolorcm{green}{blue}{red}|\\
% |\rulerparams{1pt}{}{}{5mm}{}|\\
% |\fgrulernoborderline|\\
% |\ruler{rightdown}{3cm}}|\\[2mm]
% {\fgrulernoborderline
% \fgrulercolorcm{green}{blue}{red}
% \rulerparams{1pt}{}{}{5mm}{}
% \ruler{rightdown}{3cm}}
%
% \item\commandinline|\fgrulercolorin{!<颜色1>!}{!<颜色2>!}{!<颜色3>!}{!<颜色4>!}{!<颜色5>!}|\\
% 英制标尺中的刻度颜色。
% 如果某个参数为空，则该参数不被修改。
% \begin{description}
% \item\commandinline|!<颜色1>!|
% $k/16$\,inch 处的刻度颜色，其中 $k$ 为正奇数。
%
% \item\commandinline|!<颜色2>!|
% $k/8$\,inch 处的刻度颜色，其中 $k$ 为正奇数。
%
% \item\commandinline|!<颜色3>!|
% $k/4$\,inch 处的刻度颜色，其中 $k$ 为正奇数。
%
% \item\commandinline|!<颜色4>!|
% $k/2$\,inch 处的刻度颜色，其中 $k$ 为正奇数。
%
% \item\commandinline|!<颜色5>!|
% 整数单位处的刻度颜色。
% \end{description}
%
% 默认值分别由 |\rulerparams| 的 \commandinline|!<颜色>!| 参数和 |color| 选项给出。
%
% 示例：\\
% |{\fgrulercolorin{yellow}{orange}{green}{blue}{red}|\\
% |\rulerparams{1pt}{}{}{5mm}{}|\\
% |\fgrulernoborderline|\\
% |\ruler[in]{rightdown}{3in}}|\\[2mm]
% {\fgrulernoborderline
% \fgrulercolorin{yellow}{orange}{green}{blue}{red}
% \rulerparams{1pt}{}{}{5mm}{}
% \ruler[in]{rightdown}{3in}}
%
% \item\commandinline|\fgrulerreset|\\
% 将所有选项和参数恢复为默认值。
% 此命令仅可在正文中展开。
% \end{description}
%
% \medskip\noindent
% \colorbox{red!10}{\parbox{\dimexpr\textwidth-2\fboxsep\relax}{\textcolor{red}{\faWarning}
% \itshape 所有设置命令均遵循正常的作用域规则，即如果在分组内使用，参数的更改在分组外无效。}}
% 
% \newpage\fgrulerreset
% \section{示例}
% \subsection{默认情况}\label{example:upperleft}
% \begin{examplelst}
% \documentclass{article}
% \usepackage{fgruler}
% \begin{document}
% ...
% \end{document}
% \end{examplelst}
% 
% \newpage\fgrulerreset
% \subsection{\texttt{showframe} 与 \texttt{color} 选项}\label{example:showframe-color}
% \setfgruler{color=red,showframe}
% \begin{examplelst}
% \documentclass{article}
% \usepackage[color=red,showframe]{fgruler}
% \begin{document}
% ...
% \end{document}
% \end{examplelst}
% 
% \newpage\fgrulerreset
% \subsection{\texttt{type=upperright} 选项}\label{example:upperright}
% \setfgruler{type=upperright}
% \begin{examplelst}
% \documentclass{article}
% \usepackage[type=upperright]{fgruler}
% \begin{document}
% ...
% \end{document}
% \end{examplelst}
% 
% \newpage\fgrulerreset
% \subsection{\texttt{type=lowerleft} 选项}\label{example:lowerleft}
% \setfgruler{type=lowerleft}
% \begin{examplelst}
% \documentclass{article}
% \usepackage[type=lowerleft]{fgruler}
% \begin{document}
% ...
% \end{document}
% \end{examplelst}
% 
% \newpage\fgrulerreset
% \subsection{\texttt{type=lowerright} 选项}\label{example:lowerright}
% \setfgruler{type=lowerright}
% \begin{examplelst}
% \documentclass{article}
% \usepackage[type=lowerright]{fgruler}
% \begin{document}
% ...
% \end{document}
% \end{examplelst}
% 
% \newpage\fgrulerreset
% \subsection{默认情况的偏移}\label{example:shift-first}
% \setfgruler{hshift=1cm,vshift=2cm}
% \begin{examplelst}
% \documentclass{article}
% \usepackage[hshift=1cm,vshift=2cm]{fgruler}
% \begin{document}
% ...
% \end{document}
% \end{examplelst}
% 
% \newpage\fgrulerreset
% \subsection{\texttt{type=upperright} 选项时的偏移}
% \setfgruler{type=upperright,hshift=1cm,vshift=2cm}
% \begin{examplelst}
% \documentclass{article}
% \usepackage[type=upperright,hshift=1cm,vshift=2cm]{fgruler}
% \begin{document}
% ...
% \end{document}
% \end{examplelst}
% 
% \newpage\fgrulerreset
% \subsection{\texttt{type=lowerleft} 选项时的偏移}
% \setfgruler{type=lowerleft,hshift=1cm,vshift=2cm}
% \begin{examplelst}
% \documentclass{article}
% \usepackage[type=lowerleft,hshift=1cm,vshift=2cm]{fgruler}
% \begin{document}
% ...
% \end{document}
% \end{examplelst}
% 
% \newpage\fgrulerreset
% \subsection{\texttt{type=lowerright} 选项时的偏移}\label{example:shift-last}
% \setfgruler{type=lowerright,hshift=1cm,vshift=2cm}
% \begin{examplelst}
% \documentclass{article}
% \usepackage[type=lowerright,hshift=1cm,vshift=2cm]{fgruler}
% \begin{document}
% ...
% \end{document}
% \end{examplelst}
% 
% \newpage\fgrulerreset
% \subsection{\texttt{type=alledges} 选项}\label{example:alledges}
% \setfgruler{type=alledges}
% \begin{examplelst}
% \documentclass{article}
% \usepackage[type=alledges]{fgruler}
% \begin{document}
% ...
% \end{document}
% \end{examplelst}
% 
% \newpage\fgrulerreset
% \subsection{\texttt{type=alledges*} 选项}\label{example:alledges*}
% \setfgruler{type=alledges*}
% \begin{examplelst}
% \documentclass{article}
% \usepackage[type=alledges*]{fgruler}
% \begin{document}
% ...
% \end{document}
% \end{examplelst}
% 
% \newpage\fgrulerreset
% \subsection{\texttt{type=upperleftT} 选项}\label{example:upperleftT}
% \setfgruler{type=upperleftT}
% \begin{examplelst}
% \documentclass{article}
% \usepackage[type=upperleftT]{fgruler}
% \begin{document}
% ...
% \end{document}
% \end{examplelst}
% 
% \newpage\fgrulerreset
% \subsection{\texttt{type=upperrightT} 选项}\label{example:upperrightT}
% \setfgruler{type=upperrightT}
% \begin{examplelst}
% \documentclass{article}
% \usepackage[type=upperrightT]{fgruler}
% \begin{document}
% ...
% \end{document}
% \end{examplelst}
% 
% \newpage\fgrulerreset
% \subsection{\texttt{type=lowerleftT} 选项}\label{example:lowerleftT}
% \setfgruler{type=lowerleftT}
% \begin{examplelst}
% \documentclass{article}
% \usepackage[type=lowerleftT]{fgruler}
% \begin{document}
% ...
% \end{document}
% \end{examplelst}
% 
% \newpage\fgrulerreset
% \subsection{\texttt{type=lowerrightT} 选项}\label{example:lowerrightT}
% \setfgruler{type=lowerrightT}
% \begin{examplelst}
% \documentclass{article}
% \usepackage[type=lowerrightT]{fgruler}
% \begin{document}
% ...
% \end{document}
% \end{examplelst}
% 
% \newpage\fgrulerreset
% \subsection{\texttt{type=alledgesT} 选项}\label{example:alledgesT}
% \setfgruler{type=alledgesT}
% \begin{examplelst}
% \documentclass{article}
% \usepackage[type=alledgesT]{fgruler}
% \begin{document}
% ...
% \end{document}
% \end{examplelst}
% 
% \newpage\fgrulerreset
% \subsection{\texttt{type=alledgesT*} 选项}\label{example:alledgesT*}
% \setfgruler{type=alledgesT*}
% \begin{examplelst}
% \documentclass{article}
% \usepackage[type=alledgesT*]{fgruler}
% \begin{document}
% ...
% \end{document}
% \end{examplelst}
%
% \newpage\fgrulerreset
% \subsection{组合类型名称}\label{example:combining}
% \setfgruler{type={alledges,alledgesT}}
% \fgrulerdefusercm{\fgrulertype{cm}{alledges}\fgrulertype{cm}{alledgesT}}
% \begin{examplelst}
% \documentclass{article}
% \usepackage[type={alledges,alledgesT}]{fgruler}
% \begin{document}
% ...
% \end{document}
% \end{examplelst}
% 
% \newpage\fgrulerreset
% \subsection{设置 \texttt{type=user} 选项}\label{example:user-first}
% 在下一示例中，|type=user| 选项根据页码奇偶性激活 |type=upperright| 或 |type=upperleft|。
% \setfgruler{type=none}
% \begin{examplelst}
% \documentclass{article}
% \usepackage[type=user]{fgruler}
% \fgrulerdefuser{
%     \ifodd\value{page}\fgrulertype{\fgrulerunit}{upperright}
%         \else\fgrulertype{\fgrulerunit}{upperleft}\fi
% }
% \begin{document}
% ...
% \end{document}
% \end{examplelst}
% 
% \newpage\fgrulerreset
% \subsection{设置 \texttt{type=user} 选项}
% 在下一示例中，|type=user| 选项在文本区域左边界放置一个垂直标尺。
% \setfgruler{type=user}
%\fgrulerdefuser{\AtTextLowerLeft{\ignorespaces\rulerparamsfromfg\llap{\ruler[\fgrulerunit]{downleft}{\textheight}}}}
% \begin{examplelst}
% \documentclass{article}
% \usepackage[type=user]{fgruler}
% \fgrulerdefuser{
%     \AtTextLowerLeft{% 参见 eso-pic 宏包！
%         \rulerparamsfromfg
%         \llap{\ruler[\fgrulerunit]{downleft}{\textheight}}
%     }
% }
% \begin{document}
% ...
% \end{document}
% \end{examplelst}
% 
% \newpage\fgrulerreset
% \subsection{设置 \texttt{type=user} 选项}
% 在下一示例中，|type=user| 选项在文本区域右边界和下边界放置标尺。
% \setfgruler{type=user}
%\fgrulerdefuser{\AtTextLowerLeft{\ignorespaces\rulerparamsfromfg\rulernorotatenum\llap{\ruler[\fgrulerunit]{downleft}{\textheight}}\ruler*[\fgrulerunit]{rightdown}{\textwidth}}}
% \begin{examplelst}
% \documentclass{article}
% \usepackage[type=user]{fgruler}
% \fgrulerdefuser{
%     \AtTextLowerLeft{% 参见 eso-pic 宏包！
%         \rulerparamsfromfg
%         \rulernorotatenum
%         \llap{\ruler[\fgrulerunit]{downleft}{\textheight}}%
%         \ruler*[\fgrulerunit]{rightdown}{\textwidth}
%     }
% }
% \begin{document}
% ...
% \end{document}
% \end{examplelst}
% 
% \newpage\fgrulerreset
% \subsection{设置 \texttt{type=user} 选项}\label{example:user-last}
% 在下一示例中，|type=user| 选项在文本区域右边界和上边界放置标尺。
% \setfgruler{type=user}
%\fgrulerdefuser{\AtTextUpperLeft{\ignorespaces%
%\rulerparamsfromfg\ruler[\fgrulerunit]{rightup}{\textwidth}\ignorespaces%
%\rulernorotatenum\fgrulercaptioncm{}\ignorespaces\ruler*[\fgrulerunit]{downright}{\textheight}}}
% \begin{examplelst}
% \documentclass{article}
% \usepackage[type=user]{fgruler}
% \fgrulerdefuser{
%     \AtTextUpperLeft{% 参见 eso-pic 宏包！
%         \rulerparamsfromfg
%         \ruler[\fgrulerunit]{rightup}{\textwidth}%
%         \rulernorotatenum
%         \csname fgrulercaption\fgrulerunit\endcsname{}%
%         \ruler*[\fgrulerunit]{downright}{\textheight}
%     }
% }
% \begin{document}
% ...
% \end{document}
% \end{examplelst}
%
% \newpage\fgrulerreset
% \subsection{当前页前景层与正文中的标尺}
% \setfgruler{color=blue}
% \fgruler{upperleft}{1cm}{1.5cm}\fgruler*[color=red,unit=in]{lowerrightT}
% \noindent
% text
% \rulerparams{}{\color{red}\tiny\ttfamily}{green}{}{}
% {\fgrulernoborderline\ruler{rightdown}{3cm}}
% text
% \ruler*{rightdown}{3cm}
% text
% \rotatebox[origin=tl]{30}{\ruler*{rightdown}{3cm}}
% 
% \bigskip
% \begin{examplelst}
% \documentclass{article}
% \usepackage[color=blue]{fgruler}
% \begin{document}
%     \fgruler{upperleft}{1cm}{1.5cm}
%     \fgruler*[color=red,unit=in]{lowerrightT}
%     \noindent
%     text
%     \rulerparams{}{\color{red}\tiny\ttfamily}{green}{}{}
%     {\fgrulernoborderline\ruler{rightdown}{3cm}}
%     text
%     \ruler*{rightdown}{3cm}
%     text
%     \rotatebox[origin=tl]{30}{\ruler*{rightdown}{3cm}} 
%     % \rotatebox 定义于 graphicx 宏包
% \end{document}
% \end{examplelst}
% 
% \newpage\fgrulerreset
% \subsection{正文中的标尺类型}
% \setfgruler{type=none}
% \noindent
% \rulerparams{}{}{red}{}{3pt}
% \ruler{rightdown}{3cm}
% \hfill
% \ruler{rightup}{3cm}
% \hfill
% \ruler{leftup}{3cm}
% \hfill
% \ruler{leftdown}{3cm}
% 
% \bigskip\noindent
% \rulerparams{}{}{green}{}{}
% {\rulernorotatenum\ruler{upright}{3cm}}
% \hfill
% \ruler{downright}{3cm}
% \hfill
% \ruler{upleft}{3cm}
% \hfill
% \ruler{downleft}{3cm}
% 
% \bigskip\noindent
% \rulerparams{}{}{blue!50!black}{}{}
% {\rulernorotatenum\fgrulercaptioncm{}\squareruler{upperleft}{2cm}{3cm}}
% \hfill
% \squareruler{lowerright}{2cm}{3cm}
% \hfill
% \squareruler{lowerleft}{2cm}{3cm}
% \hfill
% \squareruler{upperright}{2cm}{3cm}
% \hfill
% {\rulerparams{}{\footnotesize\bfseries\color{red}}{}{5mm}{}\squareruler[in]{lowerleft}{2in}{3cm}}
% 
% \bigskip
% \begin{examplelst}
% \documentclass{article}
% \usepackage[a4paper,margin=25mm]{geometry}
% \usepackage[nonefgrulers]{fgruler}
% \begin{document}
%     \noindent
%     \rulerparams{}{}{red}{}{3pt}
%     \ruler{rightdown}{3cm}
%     \hfill
%     \ruler{rightup}{3cm}
%     \hfill
%     \ruler{leftup}{3cm}
%     \hfill
%     \ruler{leftdown}{3cm}
%     
%     \bigskip\noindent
%     \rulerparams{}{}{green}{}{}
%     {\rulernorotatenum\ruler{upright}{3cm}}
%     \hfill
%     \ruler{downright}{3cm}
%     \hfill
%     \ruler{upleft}{3cm}
%     \hfill
%     \ruler{downleft}{3cm}
%     
%     \bigskip\noindent
%     \rulerparams{}{}{blue!50!black}{}{}
%     {\rulernorotatenum\fgrulercaptioncm{}\squareruler{upperleft}{2cm}{3cm}}
%     \hfill
%     \squareruler{lowerright}{2cm}{3cm}
%     \hfill
%     \squareruler{lowerleft}{2cm}{3cm}
%     \hfill
%     \squareruler{upperright}{2cm}{3cm}
%     \hfill
%     {\rulerparams{}{\footnotesize\bfseries\color{red}}{}{5mm}{}
%     \squareruler[in]{lowerleft}{2in}{3cm}}
% \end{document}
% \end{examplelst}
% 
% \newpage\fgrulerreset
% \subsection{卷尺}\label{example:tape}
% \setfgruler{type=none}
% 
%\noindent\ignorespaces\ruler{taperight}{\textwidth}\\[2pt]\rotatebox[origin=br]{-90}{\ruler{tapeleft}{3cm}}~\ruler{taperight}{10cm}
%\par\medskip\noindent{\fgrulerdefnum{}\ruler{taperight}{\textwidth}}
%\par\medskip\noindent{\fgrulerdefnum{}\fgrulercaptioncm{}\rulerparams{}{}{red}{}{0pt}\ruler{taperight}{\textwidth}}
%
% \bigskip
% \begin{examplelst}
% \documentclass{article}
% \usepackage[a4paper,margin=25mm]{geometry}
% \usepackage[nonefgrulers]{fgruler}
% \begin{document}
%     \noindent
%     \ruler{taperight}{\textwidth}\\[2pt]
%     \rotatebox[origin=br]{-90}{\ruler{tapeleft}{3cm}}
%     \ruler{taperight}{10cm}
%
%     \medskip\noindent
%     {\fgrulerdefnum{}
%     \ruler{taperight}{\textwidth}}
%     
%     \medskip\noindent
%     {\fgrulerdefnum{}
%     \fgrulercaptioncm{}
%     \rulerparams{}{}{red}{}{0pt}
%     \ruler{taperight}{\textwidth}}
% \end{document}
% \end{examplelst}
% 
% \newpage\fgrulerreset
% \subsection{刻度长度与旋转}
% \setfgruler{type=none}
% \noindent
% {\fgrulerdefnum{\rotatebox{45}{\arabic{fgrulernum}\,cm}}
% \fgrulercaptioncm{}
% \rulerparams{}{\tiny\color{red}}{blue}{8mm}{}
% \fgrulercolorcm{}{}{black}
% \rotatebox{-45}{\ruler{rightup}{10cm}}\\
% \ruler{rightup}{5cm}}
% 
% \bigskip
% \begin{examplelst}
% \documentclass{article}
% \usepackage[nonefgrulers]{fgruler}
% \begin{document}
%     \noindent
%     {\fgrulerdefnum{\rotatebox{45}{\arabic{fgrulernum}\,cm}}
%     \fgrulercaptioncm{}
%     \rulerparams{}{\tiny\color{red}}{blue}{8mm}{}
%     \fgrulercolorcm{}{}{black}
%     \rotatebox{-45}{\ruler{rightup}{10cm}}\\
%     \ruler{rightup}{5cm}}
% \end{document}
% \end{examplelst}
% 
% \newpage\fgrulerreset
% \subsection{坐标系}
% \setfgruler{type=none}
% \noindent
% \rulernorotatenum
% \fgrulercaptioncm{}
% \fgrulercolorcm{}{}{red}
% \rulerparams{}{\scriptsize\color{red}}{}{}{}
% {\fgrulerdefnum{$-\arabic{fgrulernum}$}\squareruler*{upperright}{3cm}{3cm}}\ignorespaces
% \squareruler{lowerleft}{13cm}{6cm}
% 
% \bigskip
% \begin{examplelst}
% \documentclass{article}
% \usepackage[a4paper,margin=25mm]{geometry}
% \usepackage[nonefgrulers]{fgruler}
% \begin{document}
%     \noindent
%     \rulernorotatenum
%     \fgrulercaptioncm{}
%     \fgrulercolorcm{}{}{red}
%     \rulerparams{}{\scriptsize\color{red}}{}{}{}
%     {\fgrulerdefnum{$-\arabic{fgrulernum}$}\squareruler*{upperright}{3cm}{3cm}}%
%     \squareruler{lowerleft}{13cm}{6cm}
% \end{document}
% \end{examplelst}
%
% \newpage\fgrulerreset
% \subsection{一种新的方形标尺类型}
% \setfgruler{type=none}
% {\rulernorotatenum\fgrulercaptioncm{}\ruler*{downleft}{5cm}}\ruler{rightup}{3cm}
% 
% \bigskip
% \begin{examplelst}
% \documentclass{article}
% \usepackage[type=none]{fgruler}
% \newcommand{\usersquareruler}[2]{%
%     {\rulernorotatenum\fgrulercaptioncm{}\ruler*{downleft}{#2}}%
%     \ruler{rightup}{#1}%
% }
% \begin{document}
% \usersquareruler{3cm}{5cm}
% \end{document}
% \end{examplelst}
%
% \newpage\fgrulerreset
% \subsection{地图上的比例尺}
% \setfgruler{type=none}
% \def\Sstackgap{5pt}
% \fgrulerratiocm{0}{0}
% \fgrulercaptioncm{m}
% \fgrulerdefnum{\fpeval{\value{fgrulernum}*100}}
% \Shortstack{\ruler{rightup}{5cm} $1:10000$}
% \par\bigskip\noindent
% \fgrulerdefnum{\fpeval{\value{fgrulernum}*200}}
% \Shortstack{\ruler{rightup}{5cm} $1:20000$}
% \par\bigskip\noindent
% \fgrulercaptioncm{km}
% \fgrulerdefnum{\fpeval{\value{fgrulernum}*.5}}
% \Shortstack{\ruler{rightup}{5cm} $1:50000$}
% 
% \bigskip
% \begin{examplelst}
% \documentclass{article}
% \usepackage[nonefgrulers]{fgruler}
% \usepackage{stackengine}
% \def\Sstackgap{5pt}
% \fgrulerratiocm{0}{0}
% \fgrulercaptioncm{m}

% \begin{document}

% \fgrulerdefnum{\fpeval{\value{fgrulernum}*100}}
% \Shortstack{\ruler{rightup}{5cm} $1:10000$}

% \bigskip

% \fgrulerdefnum{\fpeval{\value{fgrulernum}*200}}
% \Shortstack{\ruler{rightup}{5cm} $1:20000$}

% \bigskip

% \fgrulercaptioncm{km}
% \fgrulerdefnum{\fpeval{\value{fgrulernum}*.5}}
% \Shortstack{\ruler{rightup}{5cm} $1:50000$}

% \end{document}
% \end{examplelst}
%
% \StopEventually
%
%    \begin{macrocode}
%%
\@ifpackageloaded{kvoptions}{}{\RequirePackage{kvoptions}}
\@ifpackageloaded{etoolbox}{}{\RequirePackage{etoolbox}}

\SetupKeyvalOptions{family=fgruler,prefix=fgruler@}
\DeclareStringOption[.4pt]{markthick}
\DeclareStringOption[\scriptsize\sffamily]{numfont}
\DeclareStringOption[black]{color}
\DeclareStringOption[2mm]{marklength}
\DeclareStringOption[3pt]{numsep}
\DeclareStringOption[0pt]{hshift}
\DeclareStringOption[0pt]{vshift}
\DeclareStringOption[upperleft]{type}
\DeclareStringOption[cm]{unit}
\DeclareBoolOption{nonefgrulers}
\DeclareBoolOption{showframe}
\ProcessKeyvalOptions{fgruler}

\newlength{\fgruler@marklth}
\newlength{\fgruler@sep}
\newlength{\fgruler@width}
\newlength{\fgruler@fg@width}
\newlength{\fgruler@fg@height}

\def\fgruler@set#1#2#3#4#5{%
\def\fgruler@markthickness{\linethickness{#1}}%
\def\fgruler@font{#2}%
\def\fgruler@rulercolor{\color{#3}}%
\setlength{\fgruler@marklth}{#4}%
\setlength{\fgruler@sep}{#5}%
\def\fgruler@font@{\normalfont\normalsize\fgruler@font}%
\addtolength{\fgruler@sep}{\fgruler@marklth}%
\ifx\thefgrulernum\@empty%
  \expandafter\ifx\csname fgruler@caption@\fgruler@unit\endcsname\@empty%
    \setlength{\fgruler@width}{\fgruler@marklth}%
  \else
    \settoheight{\fgruler@width}{\fgruler@font@\csname fgruler@caption@\fgruler@unit\endcsname}%
    \addtolength{\fgruler@width}{\fgruler@sep}%
  \fi
\else
  \settoheight{\fgruler@width}{\fgruler@font@\thefgrulernum}%
  \addtolength{\fgruler@width}{\fgruler@sep}%
\fi}

\def\rulerparams#1#2#3#4#5{%
\ifx#1\@empty\else\def\fgruler@markthick@{#1}\fi%
\ifx#2\@empty\else\def\fgruler@numfont@{#2}\fi%
\ifx#3\@empty\else\def\fgruler@color@{#3}\fi%
\ifx#4\@empty\else\def\fgruler@marklength@{#4}\fi%
\ifx#5\@empty\else\def\fgruler@numsep@{#5}\fi%
\ignorespaces}

\rulerparams{.4pt}{\scriptsize\sffamily}{black}{2mm}{3pt}

\def\rulerparamsfromfg{\rulerparams{\fgruler@markthick}{\fgruler@numfont}{\fgruler@color}{\fgruler@marklength}{\fgruler@numsep}}

\def\fgruler@fgsetting{%
\fgruler@set{\fgruler@markthick}{\fgruler@numfont}{\fgruler@color}{\fgruler@marklength}{\fgruler@numsep}%
\rulernorotatenum}

\def\fgruler@intextsetting{%
\fgruler@set{\fgruler@markthick@}{\fgruler@numfont@}{\fgruler@color@}{\fgruler@marklength@}{\fgruler@numsep@}}

\def\fgruler@activate@type{%
\def\fgruler@output{}%
\def\fgruler@check{cm}\ifx\fgruler@unit\fgruler@check\else%
\def\fgruler@check{in}\ifx\fgruler@unit\fgruler@check\else%
\PackageError{fgruler}{\fgruler@unit\space is invalid unit}{Valid units: cm, in}%
\fi\fi%
\def\do##1{%
\expandafter\def\expandafter\fgruler@type@item\expandafter{##1}%
\def\fgruler@check{upperleft}\ifx\fgruler@type@item\fgruler@check\else%
\def\fgruler@check{upperright}\ifx\fgruler@type@item\fgruler@check\else%
\def\fgruler@check{lowerleft}\ifx\fgruler@type@item\fgruler@check\else%
\def\fgruler@check{lowerright}\ifx\fgruler@type@item\fgruler@check\else%
\def\fgruler@check{upperleftT}\ifx\fgruler@type@item\fgruler@check\else%
\def\fgruler@check{upperrightT}\ifx\fgruler@type@item\fgruler@check\else%
\def\fgruler@check{lowerleftT}\ifx\fgruler@type@item\fgruler@check\else%
\def\fgruler@check{lowerrightT}\ifx\fgruler@type@item\fgruler@check\else%
\def\fgruler@check{alledges}\ifx\fgruler@type@item\fgruler@check\else%
\def\fgruler@check{alledges*}\ifx\fgruler@type@item\fgruler@check\else%
\def\fgruler@check{alledgesT}\ifx\fgruler@type@item\fgruler@check\else%
\def\fgruler@check{alledgesT*}\ifx\fgruler@type@item\fgruler@check\else%
\def\fgruler@check{none}\ifx\fgruler@type@item\fgruler@check\else%
\def\fgruler@check{user}\ifx\fgruler@type@item\fgruler@check\else%
\PackageError{fgruler}{\fgruler@type@item\space is invalid type}%
{Valid types:\MessageBreak
upperleft, upperright, lowerleft, lowerright,\MessageBreak
upperleftT, upperrightT, lowerleftT, lowerrightT,\MessageBreak
alledges, alledges*, alledgesT, alledgesT*, user, none}%
\fi\fi\fi\fi\fi\fi\fi\fi\fi\fi\fi\fi\fi\fi%
\g@addto@macro\fgruler@output{\csname fgruler@\fgruler@unit @##1@fg\endcsname}}%
\expandafter\docsvlist\expandafter{\fgruler@type}}

\def\fgrulerratiocm#1#2{%
\ifx#1\@empty\else\def\fgruler@cm@ratio@i{#1}\fi%
\ifx#2\@empty\else\def\fgruler@cm@ratio@ii{#2}\fi%
\ignorespaces}

\fgrulerratiocm{.5}{.75}

\def\fgrulerratioin#1#2#3#4{%
\ifx#1\@empty\else\def\fgruler@in@ratio@i{#1}\fi%
\ifx#2\@empty\else\def\fgruler@in@ratio@ii{#2}\fi%
\ifx#3\@empty\else\def\fgruler@in@ratio@iii{#3}\fi%
\ifx#4\@empty\else\def\fgruler@in@ratio@iv{#4}\fi%
\ignorespaces}

\fgrulerratioin{.25}{.375}{.625}{.75}

\def\fgrulerthickcm#1#2#3{%
\ifx#1\@empty\else\def\fgruler@cm@thick@i{\linethickness{#1}}\fi%
\ifx#2\@empty\else\def\fgruler@cm@thick@ii{\linethickness{#2}}\fi%
\ifx#3\@empty\else\def\fgruler@cm@thick@iii{\linethickness{#3}}\fi%
\ignorespaces}

\def\fgruler@cm@thick@i{}
\def\fgruler@cm@thick@ii{}
\def\fgruler@cm@thick@iii{}

\def\fgrulerthickin#1#2#3#4#5{%
\ifx#1\@empty\else\def\fgruler@in@thick@i{\linethickness{#1}}\fi%
\ifx#2\@empty\else\def\fgruler@in@thick@ii{\linethickness{#2}}\fi%
\ifx#3\@empty\else\def\fgruler@in@thick@iii{\linethickness{#3}}\fi%
\ifx#4\@empty\else\def\fgruler@in@thick@iv{\linethickness{#4}}\fi%
\ifx#5\@empty\else\def\fgruler@in@thick@v{\linethickness{#5}}\fi%
\ignorespaces}

\def\fgruler@in@thick@i{}
\def\fgruler@in@thick@ii{}
\def\fgruler@in@thick@iii{}
\def\fgruler@in@thick@iv{}
\def\fgruler@in@thick@v{}

\def\fgrulercolorcm#1#2#3{%
\ifx#1\@empty\else\def\fgruler@cm@color@i{\color{#1}}\fi%
\ifx#2\@empty\else\def\fgruler@cm@color@ii{\color{#2}}\fi%
\ifx#3\@empty\else\def\fgruler@cm@color@iii{\color{#3}}\fi%
\ignorespaces}

\def\fgruler@cm@color@i{}
\def\fgruler@cm@color@ii{}
\def\fgruler@cm@color@iii{}

\def\fgrulercolorin#1#2#3#4#5{%
\ifx#1\@empty\else\def\fgruler@in@color@i{\color{#1}}\fi%
\ifx#2\@empty\else\def\fgruler@in@color@ii{\color{#2}}\fi%
\ifx#3\@empty\else\def\fgruler@in@color@iii{\color{#3}}\fi%
\ifx#4\@empty\else\def\fgruler@in@color@iv{\color{#4}}\fi%
\ifx#5\@empty\else\def\fgruler@in@color@v{\color{#5}}\fi%
\ignorespaces}

\def\fgruler@in@color@i{}
\def\fgruler@in@color@ii{}
\def\fgruler@in@color@iii{}
\def\fgruler@in@color@iv{}
\def\fgruler@in@color@v{}

\def\setfgruler#1{\setkeys{fgruler}{#1}\fgruler@activate@type\ignorespaces}

\newdimen\fgruler@tempdima
\newdimen\fgruler@tempdimb

\newcommand\fgruler@div[2]{%
\fgruler@tempdima=#1\relax\fgruler@tempdimb=\unitlength\relax
\fgruler@tempdimb=#2\fgruler@tempdimb\divide\fgruler@tempdima by \fgruler@tempdimb
\@tempcnta\fgruler@tempdima\advance\@tempcnta\@ne}

\def\fgruler@check@param#1#2{%
\expandafter\def\expandafter\fgruler@param\expandafter{#1}%
\def\fgruler@check{cm}\ifx\fgruler@param\fgruler@check\else%
\def\fgruler@check{in}\ifx\fgruler@param\fgruler@check\else%
\PackageError{fgruler}{\fgruler@param\space is invalid parameter}{Valid parameters: cm, in}%
\fi\fi%
\expandafter\def\expandafter\fgruler@param\expandafter{#2}%
\def\fgruler@check{rightdown}\ifx\fgruler@param\fgruler@check\else%
\def\fgruler@check{rightup}\ifx\fgruler@param\fgruler@check\else%
\def\fgruler@check{leftdown}\ifx\fgruler@param\fgruler@check\else%
\def\fgruler@check{leftup}\ifx\fgruler@param\fgruler@check\else%
\def\fgruler@check{downright}\ifx\fgruler@param\fgruler@check\else%
\def\fgruler@check{downleft}\ifx\fgruler@param\fgruler@check\else%
\def\fgruler@check{upright}\ifx\fgruler@param\fgruler@check\else%
\def\fgruler@check{upleft}\ifx\fgruler@param\fgruler@check\else%
\def\fgruler@check{taperight}\ifx\fgruler@param\fgruler@check\else%
\def\fgruler@check{tapeleft}\ifx\fgruler@param\fgruler@check\else%
\PackageError{fgruler}{\fgruler@param\space is invalid parameter}%
{Valid parameters:\MessageBreak
rightdown, rightup, leftdown, leftup,\MessageBreak
downright, downleft, upright, upleft, taperight, tapeleft}%
\fi\fi\fi\fi\fi\fi\fi\fi\fi\fi}

\def\fgruler@check@param@#1#2{%
\expandafter\def\expandafter\fgruler@param\expandafter{#1}%
\def\fgruler@check{cm}\ifx\fgruler@param\fgruler@check\else%
\def\fgruler@check{in}\ifx\fgruler@param\fgruler@check\else%
\PackageError{fgruler}{\fgruler@param\space is invalid parameter}{Valid parameters: cm, in}%
\fi\fi%
\expandafter\def\expandafter\fgruler@param\expandafter{#2}%
\def\fgruler@check{upperleft}\ifx\fgruler@param\fgruler@check\else%
\def\fgruler@check{upperright}\ifx\fgruler@param\fgruler@check\else%
\def\fgruler@check{lowerleft}\ifx\fgruler@param\fgruler@check\else%
\def\fgruler@check{lowerright}\ifx\fgruler@param\fgruler@check\else%
\PackageError{fgruler}{\fgruler@param\space is invalid parameter}%
{Valid parameters:\MessageBreak
upperleft, upperright, lowerleft, lowerright}%
\fi\fi\fi\fi}

\def\fgruler@check@param@@#1#2{%
\expandafter\def\expandafter\fgruler@param\expandafter{#1}%
\def\fgruler@check{cm}\ifx\fgruler@param\fgruler@check\else%
\def\fgruler@check{in}\ifx\fgruler@param\fgruler@check\else%
\PackageError{fgruler}{\fgruler@param\space is invalid parameter}{Valid parameters: cm, in}%
\fi\fi%
\expandafter\def\expandafter\fgruler@param\expandafter{#2}%
\def\fgruler@check{upperleft}\ifx\fgruler@param\fgruler@check\else%
\def\fgruler@check{upperright}\ifx\fgruler@param\fgruler@check\else%
\def\fgruler@check{lowerleft}\ifx\fgruler@param\fgruler@check\else%
\def\fgruler@check{lowerright}\ifx\fgruler@param\fgruler@check\else%
\def\fgruler@check{upperleftT}\ifx\fgruler@param\fgruler@check\else%
\def\fgruler@check{upperrightT}\ifx\fgruler@param\fgruler@check\else%
\def\fgruler@check{lowerleftT}\ifx\fgruler@param\fgruler@check\else%
\def\fgruler@check{lowerrightT}\ifx\fgruler@param\fgruler@check\else%
\def\fgruler@check{alledges}\ifx\fgruler@param\fgruler@check\else%
\def\fgruler@check{alledges*}\ifx\fgruler@param\fgruler@check\else%
\def\fgruler@check{alledgesT}\ifx\fgruler@param\fgruler@check\else%
\def\fgruler@check{alledgesT*}\ifx\fgruler@param\fgruler@check\else%
\def\fgruler@check{user}\ifx\fgruler@param\fgruler@check\else%
\PackageError{fgruler}{\fgruler@param\space is invalid type}%
{Valid types:\MessageBreak
upperleft, upperright, lowerleft, lowerright,\MessageBreak
upperleftT, upperrightT, lowerleftT, lowerrightT,\MessageBreak
alledges, alledges*, alledgesT, alledgesT*, user}%
\fi\fi\fi\fi\fi\fi\fi\fi\fi\fi\fi\fi\fi}

\newcommand{\fgruler@ruler}[3][cm]{\fgruler@check@param{#1}{#2}\fgruler@intextsetting\csname fgruler@#1@#2@\endcsname{#3}}
\newcommand{\fgruler@ruler@}[3][cm]{\fgruler@check@param{#1}{#2}\fgruler@intextsetting\csname fgruler@#1@#2@@\endcsname{#3}}
\def\ruler{\@ifstar{\fgruler@ruler@}{\fgruler@ruler}}

\newcommand{\fgruler@squareruler}[4][cm]{\fgruler@check@param@{#1}{#2}\fgruler@intextsetting\csname fgruler@#1@#2@\endcsname{#3}{#4}}
\newcommand{\fgruler@squareruler@}[4][cm]{\fgruler@check@param@{#1}{#2}\fgruler@intextsetting\csname fgruler@#1@#2@@\endcsname{#3}{#4}}
\def\squareruler{\@ifstar{\fgruler@squareruler@}{\fgruler@squareruler}}

\newcommand{\fgruler@fgruler}[4][\fgruler@unit]{\fgruler@check@param@{#1}{#2}\g@addto@macro\fgruler@output@{\csname fgruler@#1@#2@fg@\endcsname{#3}{#4}}}
\newcommand{\fgruler@fgruler@}[2][]{\g@addto@macro\fgruler@output@{{\setkeys{fgruler}{#1}\fgruler@check@param@@{\fgruler@unit}{#2}\csname fgruler@\fgruler@unit @#2@fg\endcsname}}}
\def\fgruler{\@ifstar{\fgruler@fgruler@}{\fgruler@fgruler}}

\def\fgrulercaptioncm#1{\def\fgruler@caption@cm{#1}\ignorespaces}
\def\fgruler@caption@cm{cm}
\def\fgrulercaptionin#1{\def\fgruler@caption@in{#1}\ignorespaces}
\def\fgruler@caption@in{inch}

\newif\iffgruler@borderline
\def\fgrulernoborderline{\fgruler@borderlinefalse\ignorespaces}
\def\fgrulerborderline{\fgruler@borderlinetrue\ignorespaces}
\fgruler@borderlinetrue

\newcounter{fgrulernum}
\newcounter{fgruler@check}

\def\fgruler@ifnot@divisible@five#1{%
\setcounter{fgruler@check}{\value{fgrulernum}}%
\divide\value{fgruler@check}by5\relax%
\multiply\value{fgruler@check}by5\relax%
\ifnum\value{fgruler@check}=\value{fgrulernum}\else#1\fi\stepcounter{fgrulernum}}

\def\fgruler@ifodd#1{\ifodd\value{fgrulernum}#1\fi\stepcounter{fgrulernum}}

\def\fgrulerstartnum#1{\def\fgruler@startnumh{#1}\def\fgruler@startnumv{#1}\ignorespaces}
\def\fgrulerstartnumh#1{\def\fgruler@startnumh{#1}\ignorespaces}
\def\fgrulerstartnumv#1{\def\fgruler@startnumv{#1}\ignorespaces}
\def\fgruler@startnumh{0}
\def\fgruler@startnumv{0}

\def\fgrulerdefnum#1{\def\thefgrulernum{#1}\ignorespaces}

\def\fgruler@rotatebox#1#2{#2}%

\def\rulernorotatenum{%
\ifx\@onlypreamble\@notprerr%
\def\fgruler@rotatebox##1##2{##2}%
\else\PackageError{fgruler}{Don't use \protect\rulernorotatenum\space in preamble!}{}\fi%
\ignorespaces}

\def\rulerrotatenum{%
\ifx\@onlypreamble\@notprerr%
\def\fgruler@rotatebox##1##2{\rotatebox{##1}{##2}}%
\else\PackageError{fgruler}{Don't use \protect\rulerrotatenum\space in preamble!}{}\fi%
\ignorespaces}

\def\fgrulerreset{%
\ifx\@onlypreamble\@notprerr%
\def\fgruler@cm@ratio@i{.5}%
\def\fgruler@cm@ratio@ii{.75}%
\def\fgruler@in@ratio@i{.25}%
\def\fgruler@in@ratio@ii{.375}%
\def\fgruler@in@ratio@iii{.625}%
\def\fgruler@in@ratio@iv{.75}%
\def\fgruler@cm@thick@i{}%
\def\fgruler@cm@thick@ii{}%
\def\fgruler@cm@thick@iii{}%
\def\fgruler@in@thick@i{}%
\def\fgruler@in@thick@ii{}%
\def\fgruler@in@thick@iii{}%
\def\fgruler@in@thick@iv{}%
\def\fgruler@in@thick@v{}%
\def\fgruler@cm@color@i{}%
\def\fgruler@cm@color@ii{}%
\def\fgruler@cm@color@iii{}%
\def\fgruler@in@color@i{}%
\def\fgruler@in@color@ii{}%
\def\fgruler@in@color@iii{}%
\def\fgruler@in@color@iv{}%
\def\fgruler@in@color@v{}%
\def\fgruler@rotatebox##1##2{\rotatebox{##1}{##2}}%
\def\fgruler@markthick@{.4pt}%
\def\fgruler@numfont@{\scriptsize\sffamily}%
\def\fgruler@color@{black}%
\def\fgruler@marklength@{2mm}%
\def\fgruler@numsep@{3pt}%
\def\fgruler@type{upperleft}%
\def\fgruler@unit{cm}%
\def\fgruler@markthick{.4pt}%
\def\fgruler@numfont{\scriptsize\sffamily}%
\def\fgruler@color{black}%
\def\fgruler@marklength{2mm}%
\def\fgruler@hshift{0pt}%
\def\fgruler@vshift{0pt}%
\def\fgruler@numsep{3pt}%
\def\thefgrulernum{\arabic{fgrulernum}}%
\def\fgruler@caption@cm{cm}%
\def\fgruler@caption@in{inch}%
\def\fgruler@startnumh{0}%
\def\fgruler@startnumv{0}%
\fgruler@borderlinetrue%
\fgruler@showframefalse%
\def\fgruler@output{\csname fgruler@\fgruler@unit @\fgruler@type @fg\endcsname}%
\else\PackageError{fgruler}{Don't use \protect\fgrulerreset\space in preamble!}{}\fi%
\ignorespaces}

\AtEndPreamble{

\@ifpackageloaded{xcolor}{}{\RequirePackage{xcolor}}
\@ifpackageloaded{graphicx}{}{\RequirePackage{graphicx}}

\def\fgruler@rotatebox#1#2{\rotatebox{#1}{#2}}

\iffgruler@nonefgrulers

\renewcommand{\fgruler@fgruler}[4][]{}
\renewcommand{\fgruler@fgruler@}[2][]{}

\else

\@ifpackageloaded{eso-pic}{}{\RequirePackage{eso-pic}[2020/10/14]}

\def\fgruler@frame{%
  \begingroup
    \fgruler@fgsetting
    \fgruler@rulercolor%
    \fgruler@markthickness%
    \AtPageLowerLeft{%
      \framebox(\paperwidth,\paperheight){}}%
    \AtTextLowerLeft{%
      \framebox(\textwidth,\textheight){}}%
    \AtTextUpperLeft{%
      \put(0,\headsep){%
        \framebox(\textwidth,\headheight){}}}%
    \AtTextLowerLeft{%
      \put(0,-\footskip){%
        \line(1,0){\textwidth}}}%
    \AtTextLowerLeft{%
      \fgruler@tempdima=\textwidth\advance\fgruler@tempdima\marginparsep%
      \if@twoside
        \ifodd\c@page\else
          \if@mparswitch
            \fgruler@tempdima=-\marginparsep\advance\fgruler@tempdima-\marginparwidth
          \fi
        \fi
      \fi
      \put(\fgruler@tempdima,0)%
        {\framebox(\marginparwidth,\textheight){}}%
    }%
  \endgroup
}

\def\fgruler@output@{}

\ESO@isMEMOIR{%
  \AddToHook{shipout/foreground}{%
    \ESO@MEMOIR@corr
    \put(\ESO@tempdima,\ESO@tempdimb){%
      \iffgruler@showframe\fgruler@frame\fi%
      \fgruler@output\fgruler@output@%
      \global\let\fgruler@output@\@empty%
    }%
  }%
}{%
  \AddToHook{shipout/foreground}{%
    \put(0,\ESO@yoffsetI){%
      \iffgruler@showframe\fgruler@frame\fi%
      \fgruler@output\fgruler@output@%
      \global\let\fgruler@output@\@empty%
    }%
  }%
}

\fgruler@activate@type

\fi
}

%% RIGHTDOWN CM
\def\fgruler@cm@rightdown#1{%
\fgruler@rulercolor%
\fgruler@markthickness%
\fgruler@div{#1}{1}%
\setcounter{fgrulernum}{0}%
\multiput(0,0)(1,0){\@tempcnta}{\fgruler@ifnot@divisible@five{\fgruler@cm@thick@i\fgruler@cm@color@i\line(0,-1){\fgruler@cm@ratio@i\fgruler@marklth}}}%
\fgruler@div{#1}{5}%
\setcounter{fgrulernum}{0}%
\multiput(0,0)(5,0){\@tempcnta}{\fgruler@ifodd{\fgruler@cm@thick@ii\fgruler@cm@color@ii\line(0,-1){\fgruler@cm@ratio@ii\fgruler@marklth}}}%
\fgruler@div{#1}{10}%
\multiput(0,0)(10,0){\@tempcnta}{\fgruler@cm@thick@iii\fgruler@cm@color@iii\line(0,-1){\fgruler@marklth}}%
\iffgruler@borderline\put(-\@halfwidth,0){\line(1,0){#1+\@wholewidth}}\fi% \@wholewidth is current line thickness and \@halfwidth = 0.5\@wholewidth (see in latex.ltx)
\setcounter{fgrulernum}{\fgruler@startnumh}\stepcounter{fgrulernum}%
\multiput(10,-\fgruler@sep)(10,0){\fgruler@tempdima}{\makebox(0,0)[t]{\fgruler@font@\thefgrulernum\stepcounter{fgrulernum}}}}

\def\fgruler@cm@rightdown@#1{%
\setlength{\unitlength}{1mm}%
\begin{picture}(#1,\fgruler@width)(0,-\fgruler@width)%
\fgruler@cm@rightdown{#1}%
\setcounter{fgrulernum}{\fgruler@startnumh}%
\put(0,-\fgruler@sep){\makebox(0,0)[t]{\fgruler@font@\thefgrulernum}}%
\ifdim#1>.5cm\put(5,-\fgruler@sep){\makebox(0,0)[t]{\fgruler@font@\fgruler@caption@cm}}\fi%
\end{picture}}

\def\fgruler@cm@rightdown@@#1{\leavevmode\hbox{}\lower\fgruler@width\hbox{\fgruler@cm@rightdown@{#1}}}

%% RIGHTDOWN IN
\def\fgruler@in@rightdown#1{%
\fgruler@rulercolor%
\fgruler@markthickness%
\fgruler@div{#1}{1}%
\setcounter{fgrulernum}{0}%
\multiput(0,0)(1,0){\@tempcnta}{\fgruler@ifodd{\fgruler@in@thick@i\fgruler@in@color@i\line(0,-1){\fgruler@in@ratio@i\fgruler@marklth}}}%
\fgruler@div{#1}{2}%
\setcounter{fgrulernum}{0}%
\multiput(0,0)(2,0){\@tempcnta}{\fgruler@ifodd{\fgruler@in@thick@ii\fgruler@in@color@ii\line(0,-1){\fgruler@in@ratio@ii\fgruler@marklth}}}%
\fgruler@div{#1}{4}%
\setcounter{fgrulernum}{0}%
\multiput(0,0)(4,0){\@tempcnta}{\fgruler@ifodd{\fgruler@in@thick@iii\fgruler@in@color@iii\line(0,-1){\fgruler@in@ratio@iii\fgruler@marklth}}}%
\fgruler@div{#1}{8}%
\setcounter{fgrulernum}{0}%
\multiput(0,0)(8,0){\@tempcnta}{\fgruler@ifodd{\fgruler@in@thick@iv\fgruler@in@color@iv\line(0,-1){\fgruler@in@ratio@iv\fgruler@marklth}}}%
\fgruler@div{#1}{16}%
\multiput(0,0)(16,0){\@tempcnta}{\fgruler@in@thick@v\fgruler@in@color@v\line(0,-1){\fgruler@marklth}}%
\iffgruler@borderline\put(-\@halfwidth,0){\line(1,0){#1+\@wholewidth}}\fi%
\setcounter{fgrulernum}{\fgruler@startnumh}\stepcounter{fgrulernum}%
\multiput(16,-\fgruler@sep)(16,0){\fgruler@tempdima}{\makebox(0,0)[t]{\fgruler@font@\thefgrulernum\stepcounter{fgrulernum}}}}

\def\fgruler@in@rightdown@#1{%
\setlength{\unitlength}{.0625in}%
\begin{picture}(#1,\fgruler@width)(0,-\fgruler@width)%
\fgruler@in@rightdown{#1}%
\setcounter{fgrulernum}{\fgruler@startnumh}%
\put(0,-\fgruler@sep){\makebox(0,0)[t]{\fgruler@font@\thefgrulernum}}%
\ifdim#1>.5in\put(8,-\fgruler@sep){\makebox(0,0)[t]{\fgruler@font@\fgruler@caption@in}}\fi%
\end{picture}}

\def\fgruler@in@rightdown@@#1{\leavevmode\hbox{}\lower\fgruler@width\hbox{\fgruler@in@rightdown@{#1}}}

%% RIGHTUP CM
\def\fgruler@cm@rightup#1{%
\fgruler@rulercolor%
\fgruler@markthickness%
\fgruler@div{#1}{1}%
\setcounter{fgrulernum}{0}%
\multiput(0,0)(1,0){\@tempcnta}{\fgruler@ifnot@divisible@five{\fgruler@cm@thick@i\fgruler@cm@color@i\line(0,1){\fgruler@cm@ratio@i\fgruler@marklth}}}%
\fgruler@div{#1}{5}%
\setcounter{fgrulernum}{0}%
\multiput(0,0)(5,0){\@tempcnta}{\fgruler@ifodd{\fgruler@cm@thick@ii\fgruler@cm@color@ii\line(0,1){\fgruler@cm@ratio@ii\fgruler@marklth}}}%
\fgruler@div{#1}{10}%
\multiput(0,0)(10,0){\@tempcnta}{\fgruler@cm@thick@iii\fgruler@cm@color@iii\line(0,1){\fgruler@marklth}}%
\iffgruler@borderline\put(-\@halfwidth,0){\line(1,0){#1+\@wholewidth}}\fi%
\setcounter{fgrulernum}{\fgruler@startnumh}\stepcounter{fgrulernum}%
\multiput(10,\fgruler@sep)(10,0){\fgruler@tempdima}{\makebox(0,0)[b]{\fgruler@font@\thefgrulernum\stepcounter{fgrulernum}}}}

\def\fgruler@cm@rightup@#1{%
\setlength{\unitlength}{1mm}%
\begin{picture}(#1,\fgruler@width)%
\fgruler@cm@rightup{#1}%
\setcounter{fgrulernum}{\fgruler@startnumh}%
\put(0,\fgruler@sep){\makebox(0,0)[b]{\fgruler@font@\thefgrulernum}}%
\ifdim#1>.5cm\put(5,\fgruler@sep){\makebox(0,0)[b]{\fgruler@font@\fgruler@caption@cm}}\fi%
\end{picture}}

\def\fgruler@cm@rightup@@#1{\leavevmode\hbox{}\lower\fgruler@width\hbox{\fgruler@cm@rightup@{#1}}}

%% RIGHTUP IN
\def\fgruler@in@rightup#1{%
\fgruler@rulercolor%
\fgruler@markthickness%
\fgruler@div{#1}{1}%
\setcounter{fgrulernum}{0}%
\multiput(0,0)(1,0){\@tempcnta}{\fgruler@ifodd{\fgruler@in@thick@i\fgruler@in@color@i\line(0,1){\fgruler@in@ratio@i\fgruler@marklth}}}%
\fgruler@div{#1}{2}%
\setcounter{fgrulernum}{0}%
\multiput(0,0)(2,0){\@tempcnta}{\fgruler@ifodd{\fgruler@in@thick@ii\fgruler@in@color@ii\line(0,1){\fgruler@in@ratio@ii\fgruler@marklth}}}%
\fgruler@div{#1}{4}%
\setcounter{fgrulernum}{0}%
\multiput(0,0)(4,0){\@tempcnta}{\fgruler@ifodd{\fgruler@in@thick@iii\fgruler@in@color@iii\line(0,1){\fgruler@in@ratio@iii\fgruler@marklth}}}%
\fgruler@div{#1}{8}%
\setcounter{fgrulernum}{0}%
\multiput(0,0)(8,0){\@tempcnta}{\fgruler@ifodd{\fgruler@in@thick@iv\fgruler@in@color@iv\line(0,1){\fgruler@in@ratio@iv\fgruler@marklth}}}%
\fgruler@div{#1}{16}%
\multiput(0,0)(16,0){\@tempcnta}{\fgruler@in@thick@v\fgruler@in@color@v\line(0,1){\fgruler@marklth}}%
\iffgruler@borderline\put(-\@halfwidth,0){\line(1,0){#1+\@wholewidth}}\fi%
\setcounter{fgrulernum}{\fgruler@startnumh}\stepcounter{fgrulernum}%
\multiput(16,\fgruler@sep)(16,0){\fgruler@tempdima}{\makebox(0,0)[b]{\fgruler@font@\thefgrulernum\stepcounter{fgrulernum}}}}

\def\fgruler@in@rightup@#1{%
\setlength{\unitlength}{.0625in}%
\begin{picture}(#1,\fgruler@width)%
\fgruler@in@rightup{#1}%
\setcounter{fgrulernum}{\fgruler@startnumh}%
\put(0,\fgruler@sep){\makebox(0,0)[b]{\fgruler@font@\thefgrulernum}}%
\ifdim#1>.5in\put(8,\fgruler@sep){\makebox(0,0)[b]{\fgruler@font@\fgruler@caption@in}}\fi%
\end{picture}}

\def\fgruler@in@rightup@@#1{\leavevmode\hbox{}\lower\fgruler@width\hbox{\fgruler@in@rightup@{#1}}}

%% LEFTDOWN CM
\def\fgruler@cm@leftdown#1{%
\fgruler@rulercolor%
\fgruler@markthickness%
\fgruler@div{#1}{1}%
\setcounter{fgrulernum}{0}%
\multiput(0,0)(-1,0){\@tempcnta}{\fgruler@ifnot@divisible@five{\fgruler@cm@thick@i\fgruler@cm@color@i\line(0,-1){\fgruler@cm@ratio@i\fgruler@marklth}}}%
\fgruler@div{#1}{5}%
\setcounter{fgrulernum}{0}%
\multiput(0,0)(-5,0){\@tempcnta}{\fgruler@ifodd{\fgruler@cm@thick@ii\fgruler@cm@color@ii\line(0,-1){\fgruler@cm@ratio@ii\fgruler@marklth}}}%
\fgruler@div{#1}{10}%
\multiput(0,0)(-10,0){\@tempcnta}{\fgruler@cm@thick@iii\fgruler@cm@color@iii\line(0,-1){\fgruler@marklth}}%
\iffgruler@borderline\put(\@halfwidth,0){\line(-1,0){#1+\@wholewidth}}\fi%
\setcounter{fgrulernum}{\fgruler@startnumh}\stepcounter{fgrulernum}%
\multiput(-10,-\fgruler@sep)(-10,0){\fgruler@tempdima}{\makebox(0,0)[t]{\fgruler@font@\thefgrulernum\stepcounter{fgrulernum}}}}

\def\fgruler@cm@leftdown@#1{%
\setlength{\unitlength}{1mm}%
\begin{picture}(#1,\fgruler@width)(-#1,-\fgruler@width)%
\fgruler@cm@leftdown{#1}%
\setcounter{fgrulernum}{\fgruler@startnumh}%
\put(0,-\fgruler@sep){\makebox(0,0)[t]{\fgruler@font@\thefgrulernum}}%
\ifdim#1>.5cm\put(-5,-\fgruler@sep){\makebox(0,0)[t]{\fgruler@font@\fgruler@caption@cm}}\fi%
\end{picture}}

\def\fgruler@cm@leftdown@@#1{\leavevmode\hbox{}\lower\fgruler@width\hbox{\fgruler@cm@leftdown@{#1}}}

%% LEFTDOWN IN
\def\fgruler@in@leftdown#1{%
\fgruler@rulercolor%
\fgruler@markthickness%
\fgruler@div{#1}{1}%
\setcounter{fgrulernum}{0}%
\multiput(0,0)(-1,0){\@tempcnta}{\fgruler@ifodd{\fgruler@in@thick@i\fgruler@in@color@i\line(0,-1){\fgruler@in@ratio@i\fgruler@marklth}}}%
\fgruler@div{#1}{2}%
\setcounter{fgrulernum}{0}%
\multiput(0,0)(-2,0){\@tempcnta}{\fgruler@ifodd{\fgruler@in@thick@ii\fgruler@in@color@ii\line(0,-1){\fgruler@in@ratio@ii\fgruler@marklth}}}%
\fgruler@div{#1}{4}%
\setcounter{fgrulernum}{0}%
\multiput(0,0)(-4,0){\@tempcnta}{\fgruler@ifodd{\fgruler@in@thick@iii\fgruler@in@color@iii\line(0,-1){\fgruler@in@ratio@iii\fgruler@marklth}}}%
\fgruler@div{#1}{8}%
\setcounter{fgrulernum}{0}%
\multiput(0,0)(-8,0){\@tempcnta}{\fgruler@ifodd{\fgruler@in@thick@iv\fgruler@in@color@iv\line(0,-1){\fgruler@in@ratio@iv\fgruler@marklth}}}%
\fgruler@div{#1}{16}%
\multiput(0,0)(-16,0){\@tempcnta}{\fgruler@in@thick@v\fgruler@in@color@v\line(0,-1){\fgruler@marklth}}%
\iffgruler@borderline\put(\@halfwidth,0){\line(-1,0){#1+\@wholewidth}}\fi%
\setcounter{fgrulernum}{\fgruler@startnumh}\stepcounter{fgrulernum}%
\multiput(-16,-\fgruler@sep)(-16,0){\fgruler@tempdima}{\makebox(0,0)[t]{\fgruler@font@\thefgrulernum\stepcounter{fgrulernum}}}}

\def\fgruler@in@leftdown@#1{%
\setlength{\unitlength}{.0625in}%
\begin{picture}(#1,\fgruler@width)(-#1,-\fgruler@width)%
\fgruler@in@leftdown{#1}%
\setcounter{fgrulernum}{\fgruler@startnumh}%
\put(0,-\fgruler@sep){\makebox(0,0)[t]{\fgruler@font@\thefgrulernum}}%
\ifdim#1>.5in\put(-8,-\fgruler@sep){\makebox(0,0)[t]{\fgruler@font@\fgruler@caption@in}}\fi%
\end{picture}}

\def\fgruler@in@leftdown@@#1{\leavevmode\hbox{}\lower\fgruler@width\hbox{\fgruler@in@leftdown@{#1}}}

%% LEFTUP CM
\def\fgruler@cm@leftup#1{%
\fgruler@rulercolor%
\fgruler@markthickness%
\fgruler@div{#1}{1}%
\setcounter{fgrulernum}{0}%
\multiput(0,0)(-1,0){\@tempcnta}{\fgruler@ifnot@divisible@five{\fgruler@cm@thick@i\fgruler@cm@color@i\line(0,1){\fgruler@cm@ratio@i\fgruler@marklth}}}%
\fgruler@div{#1}{5}%
\setcounter{fgrulernum}{0}%
\multiput(0,0)(-5,0){\@tempcnta}{\fgruler@ifodd{\fgruler@cm@thick@ii\fgruler@cm@color@ii\line(0,1){\fgruler@cm@ratio@ii\fgruler@marklth}}}%
\fgruler@div{#1}{10}%
\multiput(0,0)(-10,0){\@tempcnta}{\fgruler@cm@thick@iii\fgruler@cm@color@iii\line(0,1){\fgruler@marklth}}%
\iffgruler@borderline\put(\@halfwidth,0){\line(-1,0){#1+\@wholewidth}}\fi%
\setcounter{fgrulernum}{\fgruler@startnumh}\stepcounter{fgrulernum}%
\multiput(-10,\fgruler@sep)(-10,0){\fgruler@tempdima}{\makebox(0,0)[b]{\fgruler@font@\thefgrulernum\stepcounter{fgrulernum}}}}

\def\fgruler@cm@leftup@#1{%
\setlength{\unitlength}{1mm}%
\begin{picture}(#1,\fgruler@width)(-#1,0)%
\fgruler@cm@leftup{#1}%
\setcounter{fgrulernum}{\fgruler@startnumh}%
\put(0,\fgruler@sep){\makebox(0,0)[b]{\fgruler@font@\thefgrulernum}}%
\ifdim#1>.5cm\put(-5,\fgruler@sep){\makebox(0,0)[b]{\fgruler@font@\fgruler@caption@cm}}\fi%
\end{picture}}

\def\fgruler@cm@leftup@@#1{\leavevmode\hbox{}\lower\fgruler@width\hbox{\fgruler@cm@leftup@{#1}}}

%% LEFTUP IN
\def\fgruler@in@leftup#1{%
\fgruler@rulercolor%
\fgruler@markthickness%
\fgruler@div{#1}{1}%
\setcounter{fgrulernum}{0}%
\multiput(0,0)(-1,0){\@tempcnta}{\fgruler@ifodd{\fgruler@in@thick@i\fgruler@in@color@i\line(0,1){\fgruler@in@ratio@i\fgruler@marklth}}}%
\fgruler@div{#1}{2}%
\setcounter{fgrulernum}{0}%
\multiput(0,0)(-2,0){\@tempcnta}{\fgruler@ifodd{\fgruler@in@thick@ii\fgruler@in@color@ii\line(0,1){\fgruler@in@ratio@ii\fgruler@marklth}}}%
\fgruler@div{#1}{4}%
\setcounter{fgrulernum}{0}%
\multiput(0,0)(-4,0){\@tempcnta}{\fgruler@ifodd{\fgruler@in@thick@iii\fgruler@in@color@iii\line(0,1){\fgruler@in@ratio@iii\fgruler@marklth}}}%
\fgruler@div{#1}{8}%
\setcounter{fgrulernum}{0}%
\multiput(0,0)(-8,0){\@tempcnta}{\fgruler@ifodd{\fgruler@in@thick@iv\fgruler@in@color@iv\line(0,1){\fgruler@in@ratio@iv\fgruler@marklth}}}%
\fgruler@div{#1}{16}%
\multiput(0,0)(-16,0){\@tempcnta}{\fgruler@in@thick@v\fgruler@in@color@v\line(0,1){\fgruler@marklth}}%
\iffgruler@borderline\put(\@halfwidth,0){\line(-1,0){#1+\@wholewidth}}\fi%
\setcounter{fgrulernum}{\fgruler@startnumh}\stepcounter{fgrulernum}%
\multiput(-16,\fgruler@sep)(-16,0){\fgruler@tempdima}{\makebox(0,0)[b]{\fgruler@font@\thefgrulernum\stepcounter{fgrulernum}}}}

\def\fgruler@in@leftup@#1{%
\setlength{\unitlength}{.0625in}%
\begin{picture}(#1,\fgruler@width)(-#1,0)%
\fgruler@in@leftup{#1}%
\setcounter{fgrulernum}{\fgruler@startnumh}%
\put(0,\fgruler@sep){\makebox(0,0)[b]{\fgruler@font@\thefgrulernum}}%
\ifdim#1>.5in\put(-8,\fgruler@sep){\makebox(0,0)[b]{\fgruler@font@\fgruler@caption@in}}\fi%
\end{picture}}

\def\fgruler@in@leftup@@#1{\leavevmode\hbox{}\lower\fgruler@width\hbox{\fgruler@in@leftup@{#1}}}

%% DOWNRIGHT CM
\def\fgruler@cm@downright#1{%
\fgruler@rulercolor%
\fgruler@markthickness%
\fgruler@div{#1}{1}%
\setcounter{fgrulernum}{0}%
\multiput(0,0)(0,-1){\@tempcnta}{\fgruler@ifnot@divisible@five{\fgruler@cm@thick@i\fgruler@cm@color@i\line(1,0){\fgruler@cm@ratio@i\fgruler@marklth}}}%
\fgruler@div{#1}{5}%
\setcounter{fgrulernum}{0}%
\multiput(0,0)(0,-5){\@tempcnta}{\fgruler@ifodd{\fgruler@cm@thick@ii\fgruler@cm@color@ii\line(1,0){\fgruler@cm@ratio@ii\fgruler@marklth}}}%
\fgruler@div{#1}{10}%
\multiput(0,0)(0,-10){\@tempcnta}{\fgruler@cm@thick@iii\fgruler@cm@color@iii\line(1,0){\fgruler@marklth}}%
\iffgruler@borderline\put(0,\@halfwidth){\line(0,-1){#1+\@wholewidth}}\fi%
\setcounter{fgrulernum}{\fgruler@startnumv}\stepcounter{fgrulernum}%
\multiput(\fgruler@sep,-10)(0,-10){\fgruler@tempdima}{\makebox(0,0)[l]{\fgruler@rotatebox{90}{\fgruler@font@\thefgrulernum\stepcounter{fgrulernum}}}}}

\def\fgruler@cm@downright@#1{%
\setlength{\unitlength}{1mm}%
\begin{picture}(\fgruler@width,#1)(0,-#1)%
\fgruler@cm@downright{#1}%
\setcounter{fgrulernum}{\fgruler@startnumv}%
\put(\fgruler@sep,0){\makebox(0,0)[l]{\fgruler@rotatebox{90}{\fgruler@font@\thefgrulernum}}}%
\ifdim#1>.5cm\put(\fgruler@sep,-5){\makebox(0,0)[l]{\fgruler@rotatebox{90}{\fgruler@font@\fgruler@caption@cm}}}\fi%
\end{picture}}

\def\fgruler@cm@downright@@#1{\leavevmode\hbox{}\lower#1\hbox{\fgruler@cm@downright@{#1}}}

%% DOWNRIGHT IN
\def\fgruler@in@downright#1{%
\fgruler@rulercolor%
\fgruler@markthickness%
\fgruler@div{#1}{1}%
\setcounter{fgrulernum}{0}%
\multiput(0,0)(0,-1){\@tempcnta}{\fgruler@ifodd{\fgruler@in@thick@i\fgruler@in@color@i\line(1,0){\fgruler@in@ratio@i\fgruler@marklth}}}%
\fgruler@div{#1}{2}%
\setcounter{fgrulernum}{0}%
\multiput(0,0)(0,-2){\@tempcnta}{\fgruler@ifodd{\fgruler@in@thick@ii\fgruler@in@color@ii\line(1,0){\fgruler@in@ratio@ii\fgruler@marklth}}}%
\fgruler@div{#1}{4}%
\setcounter{fgrulernum}{0}%
\multiput(0,0)(0,-4){\@tempcnta}{\fgruler@ifodd{\fgruler@in@thick@iii\fgruler@in@color@iii\line(1,0){\fgruler@in@ratio@iii\fgruler@marklth}}}%
\fgruler@div{#1}{8}%
\setcounter{fgrulernum}{0}%
\multiput(0,0)(0,-8){\@tempcnta}{\fgruler@ifodd{\fgruler@in@thick@iv\fgruler@in@color@iv\line(1,0){\fgruler@in@ratio@iv\fgruler@marklth}}}%
\fgruler@div{#1}{16}%
\multiput(0,0)(0,-16){\@tempcnta}{\fgruler@in@thick@v\fgruler@in@color@v\line(1,0){\fgruler@marklth}}%
\iffgruler@borderline\put(0,\@halfwidth){\line(0,-1){#1+\@wholewidth}}\fi%
\setcounter{fgrulernum}{\fgruler@startnumv}\stepcounter{fgrulernum}%
\multiput(\fgruler@sep,-16)(0,-16){\fgruler@tempdima}{\makebox(0,0)[l]{\fgruler@rotatebox{90}{\fgruler@font@\thefgrulernum\stepcounter{fgrulernum}}}}}

\def\fgruler@in@downright@#1{%
\setlength{\unitlength}{.0625in}%
\begin{picture}(\fgruler@width,#1)(0,-#1)%
\fgruler@in@downright{#1}%
\setcounter{fgrulernum}{\fgruler@startnumv}%
\put(\fgruler@sep,0){\makebox(0,0)[l]{\fgruler@rotatebox{90}{\fgruler@font@\thefgrulernum}}}%
\ifdim#1>.5in\put(\fgruler@sep,-8){\makebox(0,0)[l]{\fgruler@rotatebox{90}{\fgruler@font@\fgruler@caption@in}}}\fi%
\end{picture}}

\def\fgruler@in@downright@@#1{\leavevmode\hbox{}\lower#1\hbox{\fgruler@in@downright@{#1}}}

%% DOWNLEFT CM
\def\fgruler@cm@downleft#1{%
\fgruler@rulercolor%
\fgruler@markthickness%
\fgruler@div{#1}{1}%
\setcounter{fgrulernum}{0}%
\multiput(0,0)(0,-1){\@tempcnta}{\fgruler@ifnot@divisible@five{\fgruler@cm@thick@i\fgruler@cm@color@i\line(-1,0){\fgruler@cm@ratio@i\fgruler@marklth}}}%
\fgruler@div{#1}{5}%
\setcounter{fgrulernum}{0}%
\multiput(0,0)(0,-5){\@tempcnta}{\fgruler@ifodd{\fgruler@cm@thick@ii\fgruler@cm@color@ii\line(-1,0){\fgruler@cm@ratio@ii\fgruler@marklth}}}%
\fgruler@div{#1}{10}%
\multiput(0,0)(0,-10){\@tempcnta}{\fgruler@cm@thick@iii\fgruler@cm@color@iii\line(-1,0){\fgruler@marklth}}%
\iffgruler@borderline\put(0,\@halfwidth){\line(0,-1){#1+\@wholewidth}}\fi%
\setcounter{fgrulernum}{\fgruler@startnumv}\stepcounter{fgrulernum}%
\multiput(-\fgruler@sep,-10)(0,-10){\fgruler@tempdima}{\makebox(0,0)[r]{\fgruler@rotatebox{-90}{\fgruler@font@\thefgrulernum\stepcounter{fgrulernum}}}}}

\def\fgruler@cm@downleft@#1{%
\setlength{\unitlength}{1mm}%
\begin{picture}(\fgruler@width,#1)(-\fgruler@width,-#1)%
\fgruler@cm@downleft{#1}%
\setcounter{fgrulernum}{\fgruler@startnumv}%
\put(-\fgruler@sep,0){\makebox(0,0)[r]{\fgruler@rotatebox{-90}{\fgruler@font@\thefgrulernum}}}%
\ifdim#1>.5cm\put(-\fgruler@sep,-5){\makebox(0,0)[r]{\fgruler@rotatebox{-90}{\fgruler@font@\fgruler@caption@cm}}}\fi%
\end{picture}}

\def\fgruler@cm@downleft@@#1{\leavevmode\hbox{}\lower#1\hbox{\fgruler@cm@downleft@{#1}}}

%% DOWNLEFT IN
\def\fgruler@in@downleft#1{%
\fgruler@rulercolor%
\fgruler@markthickness%
\fgruler@div{#1}{1}%
\setcounter{fgrulernum}{0}%
\multiput(0,0)(0,-1){\@tempcnta}{\fgruler@ifodd{\fgruler@in@thick@i\fgruler@in@color@i\line(-1,0){\fgruler@in@ratio@i\fgruler@marklth}}}%
\fgruler@div{#1}{2}%
\setcounter{fgrulernum}{0}%
\multiput(0,0)(0,-2){\@tempcnta}{\fgruler@ifodd{\fgruler@in@thick@ii\fgruler@in@color@ii\line(-1,0){\fgruler@in@ratio@ii\fgruler@marklth}}}%
\fgruler@div{#1}{4}%
\setcounter{fgrulernum}{0}%
\multiput(0,0)(0,-4){\@tempcnta}{\fgruler@ifodd{\fgruler@in@thick@iii\fgruler@in@color@iii\line(-1,0){\fgruler@in@ratio@iii\fgruler@marklth}}}%
\fgruler@div{#1}{8}%
\setcounter{fgrulernum}{0}%
\multiput(0,0)(0,-8){\@tempcnta}{\fgruler@ifodd{\fgruler@in@thick@iv\fgruler@in@color@iv\line(-1,0){\fgruler@in@ratio@iv\fgruler@marklth}}}%
\fgruler@div{#1}{16}%
\multiput(0,0)(0,-16){\@tempcnta}{\fgruler@in@thick@v\fgruler@in@color@v\line(-1,0){\fgruler@marklth}}%
\iffgruler@borderline\put(0,\@halfwidth){\line(0,-1){#1+\@wholewidth}}\fi%
\setcounter{fgrulernum}{\fgruler@startnumv}\stepcounter{fgrulernum}%
\multiput(-\fgruler@sep,-16)(0,-16){\fgruler@tempdima}{\makebox(0,0)[r]{\fgruler@rotatebox{-90}{\fgruler@font@\thefgrulernum\stepcounter{fgrulernum}}}}}

\def\fgruler@in@downleft@#1{%
\setlength{\unitlength}{.0625in}%
\begin{picture}(\fgruler@width,#1)(-\fgruler@width,-#1)%
\fgruler@in@downleft{#1}%
\setcounter{fgrulernum}{\fgruler@startnumv}%
\put(-\fgruler@sep,0){\makebox(0,0)[r]{\fgruler@rotatebox{-90}{\fgruler@font@\thefgrulernum}}}%
\ifdim#1>.5in\put(-\fgruler@sep,-8){\makebox(0,0)[r]{\fgruler@rotatebox{-90}{\fgruler@font@\fgruler@caption@in}}}\fi%
\end{picture}}

\def\fgruler@in@downleft@@#1{\leavevmode\hbox{}\lower#1\hbox{\fgruler@in@downleft@{#1}}}

%% UPRIGHT CM
\def\fgruler@cm@upright#1{%
\fgruler@rulercolor%
\fgruler@markthickness%
\fgruler@div{#1}{1}%
\setcounter{fgrulernum}{0}%
\multiput(0,0)(0,1){\@tempcnta}{\fgruler@ifnot@divisible@five{\fgruler@cm@thick@i\fgruler@cm@color@i\line(1,0){\fgruler@cm@ratio@i\fgruler@marklth}}}%
\fgruler@div{#1}{5}%
\setcounter{fgrulernum}{0}%
\multiput(0,0)(0,5){\@tempcnta}{\fgruler@ifodd{\fgruler@cm@thick@ii\fgruler@cm@color@ii\line(1,0){\fgruler@cm@ratio@ii\fgruler@marklth}}}%
\fgruler@div{#1}{10}%
\multiput(0,0)(0,10){\@tempcnta}{\fgruler@cm@thick@iii\fgruler@cm@color@iii\line(1,0){\fgruler@marklth}}%
\iffgruler@borderline\put(0,-\@halfwidth){\line(0,1){#1+\@wholewidth}}\fi%
\setcounter{fgrulernum}{\fgruler@startnumv}\stepcounter{fgrulernum}%
\multiput(\fgruler@sep,10)(0,10){\fgruler@tempdima}{\makebox(0,0)[l]{\fgruler@rotatebox{90}{\fgruler@font@\thefgrulernum\stepcounter{fgrulernum}}}}}

\def\fgruler@cm@upright@#1{%
\setlength{\unitlength}{1mm}%
\begin{picture}(\fgruler@width,#1)%
\fgruler@cm@upright{#1}%
\setcounter{fgrulernum}{\fgruler@startnumv}%
\put(\fgruler@sep,0){\makebox(0,0)[l]{\fgruler@rotatebox{90}{\fgruler@font@\thefgrulernum}}}%
\ifdim#1>.5cm\put(\fgruler@sep,5){\makebox(0,0)[l]{\fgruler@rotatebox{90}{\fgruler@font@\fgruler@caption@cm}}}\fi%
\end{picture}}

\def\fgruler@cm@upright@@#1{\leavevmode\hbox{}\lower#1\hbox{\fgruler@cm@upright@{#1}}}

%% UPRIGHT IN
\def\fgruler@in@upright#1{%
\fgruler@rulercolor%
\fgruler@markthickness%
\fgruler@div{#1}{1}%
\setcounter{fgrulernum}{0}%
\multiput(0,0)(0,1){\@tempcnta}{\fgruler@ifodd{\fgruler@in@thick@i\fgruler@in@color@i\line(1,0){\fgruler@in@ratio@i\fgruler@marklth}}}%
\fgruler@div{#1}{2}%
\setcounter{fgrulernum}{0}%
\multiput(0,0)(0,2){\@tempcnta}{\fgruler@ifodd{\fgruler@in@thick@ii\fgruler@in@color@ii\line(1,0){\fgruler@in@ratio@ii\fgruler@marklth}}}%
\fgruler@div{#1}{4}%
\setcounter{fgrulernum}{0}%
\multiput(0,0)(0,4){\@tempcnta}{\fgruler@ifodd{\fgruler@in@thick@iii\fgruler@in@color@iii\line(1,0){\fgruler@in@ratio@iii\fgruler@marklth}}}%
\fgruler@div{#1}{8}%
\setcounter{fgrulernum}{0}%
\multiput(0,0)(0,8){\@tempcnta}{\fgruler@ifodd{\fgruler@in@thick@iv\fgruler@in@color@iv\line(1,0){\fgruler@in@ratio@iv\fgruler@marklth}}}%
\fgruler@div{#1}{16}%
\multiput(0,0)(0,16){\@tempcnta}{\fgruler@in@thick@v\fgruler@in@color@v\line(1,0){\fgruler@marklth}}%
\iffgruler@borderline\put(0,-\@halfwidth){\line(0,1){#1+\@wholewidth}}\fi%
\setcounter{fgrulernum}{\fgruler@startnumv}\stepcounter{fgrulernum}%
\multiput(\fgruler@sep,16)(0,16){\fgruler@tempdima}{\makebox(0,0)[l]{\fgruler@rotatebox{90}{\fgruler@font@\thefgrulernum\stepcounter{fgrulernum}}}}}

\def\fgruler@in@upright@#1{%
\setlength{\unitlength}{.0625in}%
\begin{picture}(\fgruler@width,#1)%
\fgruler@in@upright{#1}%
\setcounter{fgrulernum}{\fgruler@startnumv}%
\put(\fgruler@sep,0){\makebox(0,0)[l]{\fgruler@rotatebox{90}{\fgruler@font@\thefgrulernum}}}%
\ifdim#1>.5in\put(\fgruler@sep,8){\makebox(0,0)[l]{\fgruler@rotatebox{90}{\fgruler@font@\fgruler@caption@in}}}\fi%
\end{picture}}

\def\fgruler@in@upright@@#1{\leavevmode\hbox{}\lower#1\hbox{\fgruler@in@upright@{#1}}}

%% UPLEFT CM
\def\fgruler@cm@upleft#1{%
\fgruler@rulercolor%
\fgruler@markthickness%
\fgruler@div{#1}{1}%
\setcounter{fgrulernum}{0}%
\multiput(0,0)(0,1){\@tempcnta}{\fgruler@ifnot@divisible@five{\fgruler@cm@thick@i\fgruler@cm@color@i\line(-1,0){\fgruler@cm@ratio@i\fgruler@marklth}}}%
\fgruler@div{#1}{5}%
\setcounter{fgrulernum}{0}%
\multiput(0,0)(0,5){\@tempcnta}{\fgruler@ifodd{\fgruler@cm@thick@ii\fgruler@cm@color@ii\line(-1,0){\fgruler@cm@ratio@ii\fgruler@marklth}}}%
\fgruler@div{#1}{10}%
\multiput(0,0)(0,10){\@tempcnta}{\fgruler@cm@thick@iii\fgruler@cm@color@iii\line(-1,0){\fgruler@marklth}}%
\iffgruler@borderline\put(0,-\@halfwidth){\line(0,1){#1+\@wholewidth}}\fi%
\setcounter{fgrulernum}{\fgruler@startnumv}\stepcounter{fgrulernum}%
\multiput(-\fgruler@sep,10)(0,10){\fgruler@tempdima}{\makebox(0,0)[r]{\fgruler@rotatebox{-90}{\fgruler@font@\thefgrulernum\stepcounter{fgrulernum}}}}}

\def\fgruler@cm@upleft@#1{%
\setlength{\unitlength}{1mm}%
\begin{picture}(\fgruler@width,#1)(-\fgruler@width,0)%
\fgruler@cm@upleft{#1}%
\setcounter{fgrulernum}{\fgruler@startnumv}%
\put(-\fgruler@sep,0){\makebox(0,0)[r]{\fgruler@rotatebox{-90}{\fgruler@font@\thefgrulernum}}}%
\ifdim#1>.5cm\put(-\fgruler@sep,5){\makebox(0,0)[r]{\fgruler@rotatebox{-90}{\fgruler@font@\fgruler@caption@cm}}}\fi%
\end{picture}}

\def\fgruler@cm@upleft@@#1{\leavevmode\hbox{}\lower#1\hbox{\fgruler@cm@upleft@{#1}}}

%% UPLEFT IN
\def\fgruler@in@upleft#1{%
\fgruler@rulercolor%
\fgruler@markthickness%
\fgruler@div{#1}{1}%
\setcounter{fgrulernum}{0}%
\multiput(0,0)(0,1){\@tempcnta}{\fgruler@ifodd{\fgruler@in@thick@i\fgruler@in@color@i\line(-1,0){\fgruler@in@ratio@i\fgruler@marklth}}}%
\fgruler@div{#1}{2}%
\setcounter{fgrulernum}{0}%
\multiput(0,0)(0,2){\@tempcnta}{\fgruler@ifodd{\fgruler@in@thick@ii\fgruler@in@color@ii\line(-1,0){\fgruler@in@ratio@ii\fgruler@marklth}}}%
\fgruler@div{#1}{4}%
\setcounter{fgrulernum}{0}%
\multiput(0,0)(0,4){\@tempcnta}{\fgruler@ifodd{\fgruler@in@thick@iii\fgruler@in@color@iii\line(-1,0){\fgruler@in@ratio@iii\fgruler@marklth}}}%
\fgruler@div{#1}{8}%
\setcounter{fgrulernum}{0}%
\multiput(0,0)(0,8){\@tempcnta}{\fgruler@ifodd{\fgruler@in@thick@iv\fgruler@in@color@iv\line(-1,0){\fgruler@in@ratio@iv\fgruler@marklth}}}%
\fgruler@div{#1}{16}%
\multiput(0,0)(0,16){\@tempcnta}{\fgruler@in@thick@v\fgruler@in@color@v\line(-1,0){\fgruler@marklth}}%
\iffgruler@borderline\put(0,-\@halfwidth){\line(0,1){#1+\@wholewidth}}\fi%
\setcounter{fgrulernum}{\fgruler@startnumv}\stepcounter{fgrulernum}%
\multiput(-\fgruler@sep,16)(0,16){\fgruler@tempdima}{\makebox(0,0)[r]{\fgruler@rotatebox{-90}{\fgruler@font@\thefgrulernum\stepcounter{fgrulernum}}}}}

\def\fgruler@in@upleft@#1{%
\setlength{\unitlength}{.0625in}%
\begin{picture}(\fgruler@width,#1)(-\fgruler@width,0)%
\fgruler@in@upleft{#1}%
\setcounter{fgrulernum}{\fgruler@startnumv}%
\put(-\fgruler@sep,0){\makebox(0,0)[r]{\fgruler@rotatebox{-90}{\fgruler@font@\thefgrulernum}}}%
\ifdim#1>.5in\put(-\fgruler@sep,8){\makebox(0,0)[r]{\fgruler@rotatebox{-90}{\fgruler@font@\fgruler@caption@in}}}\fi%
\end{picture}}

\def\fgruler@in@upleft@@#1{\leavevmode\hbox{}\lower#1\hbox{\fgruler@in@upleft@{#1}}}

%% TAPERIGHT CM
\def\fgruler@cm@taperight@#1{%
\setlength{\unitlength}{1mm}%
\begin{picture}(#1,\fgruler@width+\fgruler@sep)%
\fgruler@cm@rightup{#1}%
\setcounter{fgrulernum}{\fgruler@startnumh}%
\put(0,\fgruler@sep){\makebox(0,0)[b]{\fgruler@font@\thefgrulernum}}%
\ifdim#1>.5cm\put(5,\fgruler@sep){\makebox(0,0)[b]{\fgruler@font@\fgruler@caption@cm}}\fi%
\def\thefgrulernum{}%
\put(0,\fgruler@width+\fgruler@sep){\fgruler@cm@rightdown{#1}}%
\end{picture}}

\def\fgruler@cm@taperight@@#1{\leavevmode\hbox{}\lower\dimexpr
\fgruler@width+\fgruler@sep\relax\hbox{\fgruler@cm@taperight@{#1}}}

%% TAPERIGHT IN
\def\fgruler@in@taperight@#1{%
\setlength{\unitlength}{.0625in}%
\begin{picture}(#1,\fgruler@width+\fgruler@sep)%
\fgruler@in@rightup{#1}%
\setcounter{fgrulernum}{\fgruler@startnumh}%
\put(0,\fgruler@sep){\makebox(0,0)[b]{\fgruler@font@\thefgrulernum}}%
\ifdim#1>.5in\put(8,\fgruler@sep){\makebox(0,0)[b]{\fgruler@font@\fgruler@caption@in}}\fi%
\def\thefgrulernum{}%
\put(0,\fgruler@width+\fgruler@sep){\fgruler@in@rightdown{#1}}%
\end{picture}}

\def\fgruler@in@taperight@@#1{\leavevmode\hbox{}\lower\dimexpr
\fgruler@width+\fgruler@sep\relax\hbox{\fgruler@in@taperight@{#1}}}

%% TAPELEFT CM
\def\fgruler@cm@tapeleft@#1{%
\setlength{\unitlength}{1mm}%
\begin{picture}(#1,\fgruler@width+\fgruler@sep)(-#1,0)%
\fgruler@cm@leftup{#1}%
\setcounter{fgrulernum}{\fgruler@startnumh}%
\put(0,\fgruler@sep){\makebox(0,0)[b]{\fgruler@font@\thefgrulernum}}%
\ifdim#1>.5cm\put(-5,\fgruler@sep){\makebox(0,0)[b]{\fgruler@font@\fgruler@caption@cm}}\fi%
\def\thefgrulernum{}%
\put(0,\fgruler@width+\fgruler@sep){\fgruler@cm@leftdown{#1}}%
\end{picture}}

\def\fgruler@cm@tapeleft@@#1{\leavevmode\hbox{}\lower\dimexpr
\fgruler@width+\fgruler@sep\relax\hbox{\fgruler@cm@tapeleft@{#1}}}

%% TAPELEFT IN
\def\fgruler@in@tapeleft@#1{%
\setlength{\unitlength}{.0625in}%
\begin{picture}(#1,\fgruler@width+\fgruler@sep)(-#1,0)%
\fgruler@in@leftup{#1}%
\setcounter{fgrulernum}{\fgruler@startnumh}%
\put(0,\fgruler@sep){\makebox(0,0)[b]{\fgruler@font@\thefgrulernum}}%
\ifdim#1>.5in\put(-8,\fgruler@sep){\makebox(0,0)[b]{\fgruler@font@\fgruler@caption@in}}\fi%
\def\thefgrulernum{}%
\put(0,\fgruler@width+\fgruler@sep){\fgruler@in@leftdown{#1}}%
\end{picture}}

\def\fgruler@in@tapeleft@@#1{\leavevmode\hbox{}\lower\dimexpr
\fgruler@width+\fgruler@sep\relax\hbox{\fgruler@in@tapeleft@{#1}}}

%% UPPERLEFT CM
\def\fgruler@cm@upperleft@#1#2{%
\setlength{\unitlength}{1mm}%
\begin{picture}(#1,#2)(0,-#2)%
\fgruler@cm@rightdown{#1}%
\fgruler@cm@downright{#2}%
\ifdim#1>.5cm\put(5,-\fgruler@sep){\makebox(0,0)[t]{\fgruler@font@\fgruler@caption@cm}}\fi%
\end{picture}}

\def\fgruler@cm@upperleft@@#1#2{\leavevmode\hbox{}\lower#2\hbox{\fgruler@cm@upperleft@{#1}{#2}}}

\def\fgruler@cm@upperleft@fg@#1#2{%
\begingroup%
\setlength{\fgruler@fg@width}{\paperwidth}%
\addtolength{\fgruler@fg@width}{-#1}%
\setlength{\fgruler@fg@height}{\paperheight}%
\addtolength{\fgruler@fg@height}{-#2}%
\fgruler@fgsetting%
\AtPageLowerLeft{\put(#1,0){%
\fgruler@cm@upperleft@{\fgruler@fg@width}{\fgruler@fg@height}}}%
\endgroup}

\def\fgruler@cm@upperleft@fg{\fgruler@cm@upperleft@fg@{\fgruler@hshift}{\fgruler@vshift}}

%% UPPERLEFT IN
\def\fgruler@in@upperleft@#1#2{%
\setlength{\unitlength}{.0625in}%
\begin{picture}(#1,#2)(0,-#2)%
\fgruler@in@rightdown{#1}%
\fgruler@in@downright{#2}%
\ifdim#1>.5in\put(8,-\fgruler@sep){\makebox(0,0)[t]{\fgruler@font@\fgruler@caption@in}}\fi%
\end{picture}}

\def\fgruler@in@upperleft@@#1#2{\leavevmode\hbox{}\lower#2\hbox{\fgruler@in@upperleft@{#1}{#2}}}

\def\fgruler@in@upperleft@fg@#1#2{%
\begingroup%
\setlength{\fgruler@fg@width}{\paperwidth}%
\addtolength{\fgruler@fg@width}{-#1}%
\setlength{\fgruler@fg@height}{\paperheight}%
\addtolength{\fgruler@fg@height}{-#2}%
\fgruler@fgsetting%
\AtPageLowerLeft{\put(#1,0){%
\fgruler@in@upperleft@{\fgruler@fg@width}{\fgruler@fg@height}}}%
\endgroup}

\def\fgruler@in@upperleft@fg{\fgruler@in@upperleft@fg@{\fgruler@hshift}{\fgruler@vshift}}

%% UPPERRIGHT CM
\def\fgruler@cm@upperright@#1#2{%
\setlength{\unitlength}{1mm}%
\begin{picture}(#1,#2)(-#1,-#2)%
\fgruler@cm@leftdown{#1}%
\fgruler@cm@downleft{#2}%
\ifdim#1>.5cm\put(-5,-\fgruler@sep){\makebox(0,0)[t]{\fgruler@font@\fgruler@caption@cm}}\fi%
\end{picture}}

\def\fgruler@cm@upperright@@#1#2{\leavevmode\hbox{}\lower#2\hbox{\fgruler@cm@upperright@{#1}{#2}}}

\def\fgruler@cm@upperright@fg@#1#2{%
\begingroup%
\setlength{\fgruler@fg@width}{\paperwidth}%
\addtolength{\fgruler@fg@width}{-#1}%
\setlength{\fgruler@fg@height}{\paperheight}%
\addtolength{\fgruler@fg@height}{-#2}%
\fgruler@fgsetting%
\AtPageLowerLeft{\fgruler@cm@upperright@{\fgruler@fg@width}{\fgruler@fg@height}}%
\endgroup}

\def\fgruler@cm@upperright@fg{\fgruler@cm@upperright@fg@{\fgruler@hshift}{\fgruler@vshift}}

%% UPPERRIGHT IN
\def\fgruler@in@upperright@#1#2{%
\setlength{\unitlength}{.0625in}%
\begin{picture}(#1,#2)(-#1,-#2)%
\fgruler@in@leftdown{#1}%
\fgruler@in@downleft{#2}%
\ifdim#1>.5in\put(-8,-\fgruler@sep){\makebox(0,0)[t]{\fgruler@font@\fgruler@caption@in}}\fi%
\end{picture}}

\def\fgruler@in@upperright@@#1#2{\leavevmode\hbox{}\lower#2\hbox{\fgruler@in@upperright@{#1}{#2}}}

\def\fgruler@in@upperright@fg@#1#2{%
\begingroup%
\setlength{\fgruler@fg@width}{\paperwidth}%
\addtolength{\fgruler@fg@width}{-#1}%
\setlength{\fgruler@fg@height}{\paperheight}%
\addtolength{\fgruler@fg@height}{-#2}%
\fgruler@fgsetting%
\AtPageLowerLeft{\fgruler@in@upperright@{\fgruler@fg@width}{\fgruler@fg@height}}%
\endgroup}

\def\fgruler@in@upperright@fg{\fgruler@in@upperright@fg@{\fgruler@hshift}{\fgruler@vshift}}

%% LOWERLEFT CM
\def\fgruler@cm@lowerleft@#1#2{%
\setlength{\unitlength}{1mm}%
\begin{picture}(#1,#2)%
\fgruler@cm@rightup{#1}%
\fgruler@cm@upright{#2}%
\ifdim#1>.5cm\put(5,\fgruler@sep){\makebox(0,0)[b]{\fgruler@font@\fgruler@caption@cm}}\fi%
\end{picture}}

\def\fgruler@cm@lowerleft@@#1#2{\leavevmode\hbox{}\lower#2\hbox{\fgruler@cm@lowerleft@{#1}{#2}}}

\def\fgruler@cm@lowerleft@fg@#1#2{%
\begingroup%
\setlength{\fgruler@fg@width}{\paperwidth}%
\addtolength{\fgruler@fg@width}{-#1}%
\setlength{\fgruler@fg@height}{\paperheight}%
\addtolength{\fgruler@fg@height}{-#2}%
\fgruler@fgsetting%
\AtPageLowerLeft{\put(#1,#2){%
\fgruler@cm@lowerleft@{\fgruler@fg@width}{\fgruler@fg@height}}}%
\endgroup}

\def\fgruler@cm@lowerleft@fg{\fgruler@cm@lowerleft@fg@{\fgruler@hshift}{\fgruler@vshift}}

%% LOWERLEFT IN
\def\fgruler@in@lowerleft@#1#2{%
\setlength{\unitlength}{.0625in}%
\begin{picture}(#1,#2)%
\fgruler@in@rightup{#1}%
\fgruler@in@upright{#2}%
\ifdim#1>.5in\put(8,\fgruler@sep){\makebox(0,0)[b]{\fgruler@font@\fgruler@caption@in}}\fi%
\end{picture}}

\def\fgruler@in@lowerleft@@#1#2{\leavevmode\hbox{}\lower#2\hbox{\fgruler@in@lowerleft@{#1}{#2}}}

\def\fgruler@in@lowerleft@fg@#1#2{%
\begingroup%
\setlength{\fgruler@fg@width}{\paperwidth}%
\addtolength{\fgruler@fg@width}{-#1}%
\setlength{\fgruler@fg@height}{\paperheight}%
\addtolength{\fgruler@fg@height}{-#2}%
\fgruler@fgsetting%
\AtPageLowerLeft{\put(#1,#2){%
\fgruler@in@lowerleft@{\fgruler@fg@width}{\fgruler@fg@height}}}%
\endgroup}

\def\fgruler@in@lowerleft@fg{\fgruler@in@lowerleft@fg@{\fgruler@hshift}{\fgruler@vshift}}

%% LOWERRIGHT CM
\def\fgruler@cm@lowerright@#1#2{%
\setlength{\unitlength}{1mm}%
\begin{picture}(#1,#2)(-#1,0)%
\fgruler@cm@leftup{#1}%
\fgruler@cm@upleft{#2}%
\ifdim#1>.5cm\put(-5,\fgruler@sep){\makebox(0,0)[b]{\fgruler@font@\fgruler@caption@cm}}\fi%
\end{picture}}

\def\fgruler@cm@lowerright@@#1#2{\leavevmode\hbox{}\lower#2\hbox{\fgruler@cm@lowerright@{#1}{#2}}}

\def\fgruler@cm@lowerright@fg@#1#2{%
\begingroup%
\setlength{\fgruler@fg@width}{\paperwidth}%
\addtolength{\fgruler@fg@width}{-#1}%
\setlength{\fgruler@fg@height}{\paperheight}%
\addtolength{\fgruler@fg@height}{-#2}%
\fgruler@fgsetting%
\AtPageLowerLeft{\put(0,#2){%
\fgruler@cm@lowerright@{\fgruler@fg@width}{\fgruler@fg@height}}}%
\endgroup}

\def\fgruler@cm@lowerright@fg{\fgruler@cm@lowerright@fg@{\fgruler@hshift}{\fgruler@vshift}}

%% LOWERRIGHT IN
\def\fgruler@in@lowerright@#1#2{%
\setlength{\unitlength}{.0625in}%
\begin{picture}(#1,#2)(-#1,0)%
\fgruler@in@leftup{#1}%
\fgruler@in@upleft{#2}%
\ifdim#1>.5in\put(-8,\fgruler@sep){\makebox(0,0)[b]{\fgruler@font@\fgruler@caption@in}}\fi%
\end{picture}}

\def\fgruler@in@lowerright@@#1#2{\leavevmode\hbox{}\lower#2\hbox{\fgruler@in@lowerright@{#1}{#2}}}

\def\fgruler@in@lowerright@fg@#1#2{%
\begingroup%
\setlength{\fgruler@fg@width}{\paperwidth}%
\addtolength{\fgruler@fg@width}{-#1}%
\setlength{\fgruler@fg@height}{\paperheight}%
\addtolength{\fgruler@fg@height}{-#2}%
\fgruler@fgsetting%
\AtPageLowerLeft{\put(0,#2){%
\fgruler@in@lowerright@{\fgruler@fg@width}{\fgruler@fg@height}}}%
\endgroup}

\def\fgruler@in@lowerright@fg{\fgruler@in@lowerright@fg@{\fgruler@hshift}{\fgruler@vshift}}

%% UPPERLEFTT CM
\def\fgruler@cm@upperleftT@fg{%
\AtTextLowerLeft{%
\setlength{\unitlength}{1mm}%
\fgruler@fgsetting%
\begin{picture}(\textwidth,\textheight)(0,-\textheight)%
\fgruler@cm@rightup{\textwidth}%
\fgruler@cm@downleft{\textheight}%
\put(5,\fgruler@sep){\makebox(0,0)[b]{\fgruler@font@\fgruler@caption@cm}}%
\setcounter{fgrulernum}{\fgruler@startnumh}%
\put(0,\fgruler@sep){\makebox(0,0)[b]{\fgruler@font@\thefgrulernum}}%
\setcounter{fgrulernum}{\fgruler@startnumv}%
\put(-\fgruler@sep,0){\makebox(0,0)[r]{\fgruler@font@\thefgrulernum}}%
\end{picture}}}

%% UPPERLEFTT IN
\def\fgruler@in@upperleftT@fg{%
\AtTextLowerLeft{%
\setlength{\unitlength}{.0625in}%
\fgruler@fgsetting%
\begin{picture}(\textwidth,\textheight)(0,-\textheight)%
\fgruler@in@rightup{\textwidth}%
\fgruler@in@downleft{\textheight}%
\put(8,\fgruler@sep){\makebox(0,0)[b]{\fgruler@font@\fgruler@caption@in}}%
\setcounter{fgrulernum}{\fgruler@startnumh}%
\put(0,\fgruler@sep){\makebox(0,0)[b]{\fgruler@font@\thefgrulernum}}%
\setcounter{fgrulernum}{\fgruler@startnumv}%
\put(-\fgruler@sep,0){\makebox(0,0)[r]{\fgruler@font@\thefgrulernum}}%
\end{picture}}}

%% UPPERRIGHTT CM
\def\fgruler@cm@upperrightT@fg{%
\AtTextLowerLeft{%
\setlength{\unitlength}{1mm}%
\fgruler@fgsetting%
\begin{picture}(\textwidth,\textheight)(-\textwidth,-\textheight)%
\fgruler@cm@leftup{\textwidth}%
\fgruler@cm@downright{\textheight}%
\put(-5,\fgruler@sep){\makebox(0,0)[b]{\fgruler@font@\fgruler@caption@cm}}%
\setcounter{fgrulernum}{\fgruler@startnumh}%
\put(0,\fgruler@sep){\makebox(0,0)[b]{\fgruler@font@\thefgrulernum}}%
\setcounter{fgrulernum}{\fgruler@startnumv}%
\put(\fgruler@sep,0){\makebox(0,0)[l]{\fgruler@font@\thefgrulernum}}%
\end{picture}}}

%% UPPERRIGHTT IN
\def\fgruler@in@upperrightT@fg{%
\AtTextLowerLeft{%
\setlength{\unitlength}{.0625in}%
\fgruler@fgsetting%
\begin{picture}(\textwidth,\textheight)(-\textwidth,-\textheight)%
\fgruler@in@leftup{\textwidth}%
\fgruler@in@downright{\textheight}%
\put(-8,\fgruler@sep){\makebox(0,0)[b]{\fgruler@font@\fgruler@caption@in}}%
\setcounter{fgrulernum}{\fgruler@startnumh}%
\put(0,\fgruler@sep){\makebox(0,0)[b]{\fgruler@font@\thefgrulernum}}%
\setcounter{fgrulernum}{\fgruler@startnumv}%
\put(\fgruler@sep,0){\makebox(0,0)[l]{\fgruler@font@\thefgrulernum}}%
\end{picture}}}

%% LOWERLEFTT CM
\def\fgruler@cm@lowerleftT@fg{%
\AtTextLowerLeft{%
\setlength{\unitlength}{1mm}%
\fgruler@fgsetting%
\begin{picture}(\textwidth,\textheight)(0,0)%
\fgruler@cm@rightdown{\textwidth}%
\fgruler@cm@upleft{\textheight}%
\put(5,-\fgruler@sep){\makebox(0,0)[t]{\fgruler@font@\fgruler@caption@cm}}%
\setcounter{fgrulernum}{\fgruler@startnumh}%
\put(0,-\fgruler@sep){\makebox(0,0)[t]{\fgruler@font@\thefgrulernum}}%
\setcounter{fgrulernum}{\fgruler@startnumv}%
\put(-\fgruler@sep,0){\makebox(0,0)[r]{\fgruler@font@\thefgrulernum}}%
\end{picture}}}

%% LOWERLEFTT IN
\def\fgruler@in@lowerleftT@fg{%
\AtTextLowerLeft{%
\setlength{\unitlength}{.0625in}%
\fgruler@fgsetting%
\begin{picture}(\textwidth,\textheight)(0,0)%
\fgruler@in@rightdown{\textwidth}%
\fgruler@in@upleft{\textheight}%
\put(8,-\fgruler@sep){\makebox(0,0)[t]{\fgruler@font@\fgruler@caption@in}}%
\setcounter{fgrulernum}{\fgruler@startnumh}%
\put(0,-\fgruler@sep){\makebox(0,0)[t]{\fgruler@font@\thefgrulernum}}%
\setcounter{fgrulernum}{\fgruler@startnumv}%
\put(-\fgruler@sep,0){\makebox(0,0)[r]{\fgruler@font@\thefgrulernum}}%
\end{picture}}}

%% LOWERRIGHTT CM
\def\fgruler@cm@lowerrightT@fg{%
\AtTextLowerLeft{%
\setlength{\unitlength}{1mm}%
\fgruler@fgsetting%
\begin{picture}(\textwidth,\textheight)(-\textwidth,0)%
\fgruler@cm@leftdown{\textwidth}%
\fgruler@cm@upright{\textheight}%
\put(-5,-\fgruler@sep){\makebox(0,0)[t]{\fgruler@font@\fgruler@caption@cm}}%
\setcounter{fgrulernum}{\fgruler@startnumh}%
\put(0,-\fgruler@sep){\makebox(0,0)[t]{\fgruler@font@\thefgrulernum}}%
\setcounter{fgrulernum}{\fgruler@startnumv}%
\put(\fgruler@sep,0){\makebox(0,0)[l]{\fgruler@font@\thefgrulernum}}%
\end{picture}}}

%% LOWERRIGHTT IN
\def\fgruler@in@lowerrightT@fg{%
\AtTextLowerLeft{%
\setlength{\unitlength}{.0625in}%
\fgruler@fgsetting%
\begin{picture}(\textwidth,\textheight)(-\textwidth,0)%
\fgruler@in@leftdown{\textwidth}%
\fgruler@in@upright{\textheight}%
\put(-8,-\fgruler@sep){\makebox(0,0)[t]{\fgruler@font@\fgruler@caption@in}}%
\setcounter{fgrulernum}{\fgruler@startnumh}%
\put(0,-\fgruler@sep){\makebox(0,0)[t]{\fgruler@font@\thefgrulernum}}%
\setcounter{fgrulernum}{\fgruler@startnumv}%
\put(\fgruler@sep,0){\makebox(0,0)[l]{\fgruler@font@\thefgrulernum}}%
\end{picture}}}

%% ALLEDGES CM
\def\fgruler@cm@alledges@fg{%
\AtPageLowerLeft{%
\setlength{\unitlength}{1mm}%
\fgruler@fgsetting%
\begin{picture}(\paperwidth,\paperheight)(0,-\paperheight)%
\fgruler@cm@rightdown{\paperwidth}%
\fgruler@cm@downright{\paperheight}%
\put(0,-\paperheight){\fgruler@cm@rightup{\paperwidth}}
\put(\paperwidth,0){\fgruler@cm@downleft{\paperheight}}
\put(5,-\fgruler@sep){\makebox(0,0)[t]{\fgruler@font@\fgruler@caption@cm}}
\end{picture}}}

%% ALLEDGES IN
\def\fgruler@in@alledges@fg{%
\AtPageLowerLeft{%
\setlength{\unitlength}{.0625in}%
\fgruler@fgsetting%
\begin{picture}(\paperwidth,\paperheight)(0,-\paperheight)%
\fgruler@in@rightdown{\paperwidth}%
\fgruler@in@downright{\paperheight}%
\put(0,-\paperheight){\fgruler@in@rightup{\paperwidth}}
\put(\paperwidth,0){\fgruler@in@downleft{\paperheight}}
\put(8,-\fgruler@sep){\makebox(0,0)[t]{\fgruler@font@\fgruler@caption@in}}
\end{picture}}}

%% ALLEDGES* CM
\@namedef{fgruler@cm@alledges*@fg}{%
\AtPageLowerLeft{%
\setlength{\unitlength}{1mm}%
\fgruler@fgsetting%
\begin{picture}(\paperwidth,\paperheight)(0,-\paperheight)%
\fgruler@cm@rightdown{\paperwidth}%
\fgruler@cm@downright{\paperheight}%
\put(\paperwidth,-\paperheight){%
\fgruler@cm@leftup{\paperwidth}%
\fgruler@cm@upleft{\paperheight}}
\put(5,-\fgruler@sep){\makebox(0,0)[t]{\fgruler@font@\fgruler@caption@cm}}
\end{picture}}}

%% ALLEDGES* IN
\@namedef{fgruler@in@alledges*@fg}{%
\AtPageLowerLeft{%
\setlength{\unitlength}{.0625in}%
\fgruler@fgsetting%
\begin{picture}(\paperwidth,\paperheight)(0,-\paperheight)%
\fgruler@in@rightdown{\paperwidth}%
\fgruler@in@downright{\paperheight}%
\put(\paperwidth,-\paperheight){%
\fgruler@in@leftup{\paperwidth}%
\fgruler@in@upleft{\paperheight}}
\put(8,-\fgruler@sep){\makebox(0,0)[t]{\fgruler@font@\fgruler@caption@in}}
\end{picture}}}

%% ALLEDGEST CM
\def\fgruler@cm@alledgesT@fg{%
\AtTextLowerLeft{%
\setlength{\unitlength}{1mm}%
\fgruler@fgsetting%
\begin{picture}(\textwidth,\textheight)(0,-\textheight)%
\fgruler@cm@rightup{\textwidth}%
\fgruler@cm@downleft{\textheight}%
\put(0,-\textheight){\fgruler@cm@rightdown{\textwidth}}
\put(\textwidth,0){\fgruler@cm@downright{\textheight}}
\put(5,\fgruler@sep){\makebox(0,0)[b]{\fgruler@font@\fgruler@caption@cm}}%
\setcounter{fgrulernum}{\fgruler@startnumh}%
\put(0,\fgruler@sep){\makebox(0,0)[b]{\fgruler@font@\thefgrulernum}}%
\put(0,-\textheight-\fgruler@sep){\makebox(0,0)[t]{\fgruler@font@\thefgrulernum}}%
\setcounter{fgrulernum}{\fgruler@startnumv}%
\put(-\fgruler@sep,0){\makebox(0,0)[r]{\fgruler@font@\thefgrulernum}}%
\put(\textwidth+\fgruler@sep,0){\makebox(0,0)[l]{\fgruler@font@\thefgrulernum}}%
\end{picture}}}

%% ALLEDGEST IN
\def\fgruler@in@alledgesT@fg{%
\AtTextLowerLeft{%
\setlength{\unitlength}{.0625in}%
\fgruler@fgsetting%
\begin{picture}(\textwidth,\textheight)(0,-\textheight)%
\fgruler@in@rightup{\textwidth}%
\fgruler@in@downleft{\textheight}%
\put(0,-\textheight){\fgruler@in@rightdown{\textwidth}}
\put(\textwidth,0){\fgruler@in@downright{\textheight}}
\put(8,\fgruler@sep){\makebox(0,0)[b]{\fgruler@font@\fgruler@caption@in}}%
\setcounter{fgrulernum}{\fgruler@startnumh}%
\put(0,\fgruler@sep){\makebox(0,0)[b]{\fgruler@font@\thefgrulernum}}%
\put(0,-\textheight-\fgruler@sep){\makebox(0,0)[t]{\fgruler@font@\thefgrulernum}}%
\setcounter{fgrulernum}{\fgruler@startnumv}%
\put(-\fgruler@sep,0){\makebox(0,0)[r]{\fgruler@font@\thefgrulernum}}%
\put(\textwidth+\fgruler@sep,0){\makebox(0,0)[l]{\fgruler@font@\thefgrulernum}}%
\end{picture}}}

%% ALLEDGEST* CM
\@namedef{fgruler@cm@alledgesT*@fg}{%
\AtTextLowerLeft{%
\setlength{\unitlength}{1mm}%
\fgruler@fgsetting%
\begin{picture}(\textwidth,\textheight)(0,-\textheight)%
\fgruler@cm@rightup{\textwidth}%
\fgruler@cm@downleft{\textheight}%
\put(\textwidth,-\textheight){\fgruler@cm@leftdown{\textwidth}}
\put(\textwidth,-\textheight){\fgruler@cm@upright{\textheight}}
\put(5,\fgruler@sep){\makebox(0,0)[b]{\fgruler@font@\fgruler@caption@cm}}%
\setcounter{fgrulernum}{\fgruler@startnumh}%
\put(0,\fgruler@sep){\makebox(0,0)[b]{\fgruler@font@\thefgrulernum}}%
\put(\textwidth,-\textheight-\fgruler@sep){\makebox(0,0)[t]{\fgruler@font@\thefgrulernum}}%
\setcounter{fgrulernum}{\fgruler@startnumv}%
\put(-\fgruler@sep,0){\makebox(0,0)[r]{\fgruler@font@\thefgrulernum}}%
\put(\textwidth+\fgruler@sep,-\textheight){\makebox(0,0)[l]{\fgruler@font@\thefgrulernum}}%
\end{picture}}}

%% ALLEDGEST* IN
\@namedef{fgruler@in@alledgesT*@fg}{%
\AtTextLowerLeft{%
\setlength{\unitlength}{.0625in}%
\fgruler@fgsetting%
\begin{picture}(\textwidth,\textheight)(0,-\textheight)%
\fgruler@in@rightup{\textwidth}%
\fgruler@in@downleft{\textheight}%
\put(\textwidth,-\textheight){\fgruler@in@leftdown{\textwidth}}
\put(\textwidth,-\textheight){\fgruler@in@upright{\textheight}}
\put(8,\fgruler@sep){\makebox(0,0)[b]{\fgruler@font@\fgruler@caption@in}}%
\setcounter{fgrulernum}{\fgruler@startnumh}%
\put(0,\fgruler@sep){\makebox(0,0)[b]{\fgruler@font@\thefgrulernum}}%
\put(\textwidth,-\textheight-\fgruler@sep){\makebox(0,0)[t]{\fgruler@font@\thefgrulernum}}%
\setcounter{fgrulernum}{\fgruler@startnumv}%
\put(-\fgruler@sep,0){\makebox(0,0)[r]{\fgruler@font@\thefgrulernum}}%
\put(\textwidth+\fgruler@sep,-\textheight){\makebox(0,0)[l]{\fgruler@font@\thefgrulernum}}%
\end{picture}}}

%% USER
\def\fgruler@cm@user@fg{}
\def\fgruler@in@user@fg{}
\def\fgrulertype#1#2{\csname fgruler@#1@#2@fg\endcsname}

\def\fgrulerdefusercm#1{\def\fgruler@cm@user@fg{#1}} % for old versions
\def\fgrulerdefuserin#1{\def\fgruler@in@user@fg{#1}} % for old versions

\def\fgrulerdefuser#1{%
\expandafter\def\expandafter\fgrulerunit\expandafter{\fgruler@unit}%
\@namedef{fgruler@\fgruler@unit @user@fg}{#1}}

%    \end{macrocode}
% \Finale
\endinput